%=============================================================================
%  SmartHome PFE Report — ISET Siliana, Academic Year 2025-2026
%  Agile / Scrum methodology
%=============================================================================
\documentclass[12pt,a4paper,oneside]{report}

%--- Encoding & language -------------------------------------------------------
\usepackage[T1]{fontenc}
\usepackage[utf8]{inputenc}
\usepackage[french]{babel}
\usepackage{lmodern}

%--- Page layout ---------------------------------------------------------------
\usepackage[top=2.5cm,bottom=2.5cm,left=3cm,right=2.5cm]{geometry}
\usepackage{setspace}
\onehalfspacing
\setlength{\parskip}{0.4em}
\setlength{\parindent}{1.2em}
\setlength{\emergencystretch}{3em}

%--- Math & symbols ------------------------------------------------------------
\usepackage{amsmath,amssymb}

%--- Graphics ------------------------------------------------------------------
\usepackage{graphicx}
\usepackage{float}
\usepackage{caption}
\usepackage{subcaption}
\usepackage{tikz}
\usetikzlibrary{
  arrows.meta,positioning,shapes.geometric,
  fit,calc,backgrounds,shadows,decorations.pathreplacing,
  mindmap,trees,matrix,chains}

%--- Tables --------------------------------------------------------------------
\usepackage{booktabs}
\usepackage{longtable}
\usepackage{array}
\usepackage{tabularx}
\usepackage{multirow}
\usepackage{colortbl}
\AtBeginEnvironment{longtable}{\small}

%--- Code listings -------------------------------------------------------------
\usepackage{listings}
\usepackage{xcolor}

\definecolor{codebg}{HTML}{F8F8F8}
\definecolor{codeframe}{HTML}{CCCCCC}
\definecolor{keyword}{HTML}{0000FF}
\definecolor{comment}{HTML}{008000}
\definecolor{string}{HTML}{A31515}
\definecolor{espritblue}{HTML}{003366}
\definecolor{espritgold}{HTML}{C8A951}
\definecolor{sprintblue}{HTML}{1A5276}
\definecolor{sprintgreen}{HTML}{1E8449}
\definecolor{lightblue}{HTML}{D6EAF8}
\definecolor{lightgreen}{HTML}{D5F5E3}
\definecolor{lightyellow}{HTML}{FDFAE4}
\definecolor{lightgray}{HTML}{F2F3F4}

\lstdefinestyle{codeblock}{
  backgroundcolor=\color{codebg},
  basicstyle=\ttfamily\footnotesize,
  keywordstyle=\color{keyword}\bfseries,
  commentstyle=\color{comment}\itshape,
  stringstyle=\color{string},
  numbers=left,
  numberstyle=\tiny\color{gray},
  stepnumber=1,
  numbersep=8pt,
  frame=single,
  rulecolor=\color{codeframe},
  breaklines=true,
  breakatwhitespace=false,
  postbreak=\mbox{\textcolor{gray}{\tiny$\hookrightarrow$}\space},
  showstringspaces=false,
  tabsize=2,
  columns=fullflexible,
  xleftmargin=0.5cm,
  framexleftmargin=0.5cm
}

\lstdefinestyle{json}{
  style=codeblock,
  language=,
  morestring=[b]",
  literate=
    *{0}{{{\color{orange!70!black}0}}}{1}
     {1}{{{\color{orange!70!black}1}}}{1}
     {2}{{{\color{orange!70!black}2}}}{1}
     {3}{{{\color{orange!70!black}3}}}{1}
     {4}{{{\color{orange!70!black}4}}}{1}
     {5}{{{\color{orange!70!black}5}}}{1}
     {6}{{{\color{orange!70!black}6}}}{1}
     {7}{{{\color{orange!70!black}7}}}{1}
     {8}{{{\color{orange!70!black}8}}}{1}
     {9}{{{\color{orange!70!black}9}}}{1}
     {:}{{{\color{gray}{:}}}}{1}
     {,}{{{\color{gray}{,}}}}{1}
}

%--- JavaScript language definition for listings (not built-in) ---------------
\lstdefinelanguage{JavaScript}{
  morekeywords={break,case,catch,class,const,continue,debugger,default,
    delete,do,else,export,extends,false,finally,for,from,function,if,
    import,in,instanceof,let,new,null,of,return,static,super,switch,
    this,throw,true,try,typeof,undefined,var,void,while,with,yield,
    async,await,get,set,target,require,module},
  sensitive=true,
  morecomment=[l]{//},
  morecomment=[s]{/*}{*/},
  morestring=[b]',
  morestring=[b]",
  morestring=[b]`
}

%--- Misc utilities ------------------------------------------------------------
\usepackage{xurl}
\usepackage{microtype}
\usepackage{enumitem}
\usepackage{etoolbox}
\usepackage{fancyhdr}
\usepackage{titlesec}
\usepackage{titletoc}
\usepackage{tcolorbox}
\tcbuselibrary{skins,breakable,listingsutf8}
\usepackage{pifont}
\usepackage{rotating}

%--- Hyperlinks (load last) ----------------------------------------------------
\usepackage{hyperref}
\hypersetup{
  colorlinks  = true,
  linkcolor   = espritblue,
  urlcolor    = blue!60!black,
  citecolor   = espritblue,
  pdftitle    = {Plateforme IoT Maison Intelligente -- Rapport PFE},
  pdfauthor   = {Imen Khadhrawi, Malak Hammami},
  pdfsubject  = {Projet de Fin d'Études en Génie Informatique},
  pdfkeywords = {IoT, Maison Intelligente, MQTT, Node.js, React, MongoDB, Scrum, ESP32}
}

%--- Custom tcolorbox environments --------------------------------------------
\tcbset{
  sprintbox/.style={
    enhanced,
    colback=lightblue,
    colframe=sprintblue,
    fonttitle=\bfseries\large,
    coltitle=white,
    attach boxed title to top left={yshift=-2mm,xshift=4mm},
    boxed title style={colback=sprintblue,sharp corners},
    sharp corners=south,
    breakable
  },
  storybox/.style={
    enhanced,
    colback=lightgreen!60,
    colframe=sprintgreen,
    fonttitle=\bfseries,
    coltitle=white,
    attach boxed title to top left={yshift=-2mm,xshift=4mm},
    boxed title style={colback=sprintgreen},
    breakable,
    left=4pt,right=4pt
  },
  goalbox/.style={
    enhanced,
    colback=lightyellow,
    colframe=espritgold,
    fonttitle=\bfseries,
    coltitle=black,
    sharp corners,
    breakable
  },
  infobox/.style={
    enhanced,
    colback=lightgray,
    colframe=gray!60,
    sharp corners,
    breakable,
    left=6pt,right=6pt,top=4pt,bottom=4pt
  }
}

%--- Page headers --------------------------------------------------------------
\setlength{\headheight}{14.06094pt}
\addtolength{\topmargin}{-2.06094pt}
\pagestyle{fancy}
\fancyhf{}
\fancyhead[L]{\small\color{espritblue}\textit{SmartHome IoT Platform}}
\fancyhead[R]{\small\color{espritblue}\textit{\nouppercase{\leftmark}}}
\fancyfoot[C]{\color{espritblue}\small\thepage}
\renewcommand{\headrulewidth}{1pt}
\renewcommand{\headrule}{{\color{espritgold}\hrule width\headwidth height\headrulewidth\vskip-\headrulewidth}}

%--- Title styles (chapter / section / subsection) ----------------------------
\titleformat{\chapter}[display]
  {\normalfont\bfseries}
  {%
    \color{espritgold}\rule{\textwidth}{1.5pt}\\[3pt]%
    \large\bfseries\color{espritblue!65}\chaptertitlename~\thechapter%
  }
  {8pt}
  {\color{espritblue}\Huge\bfseries}
  [\vspace{8pt}\color{espritblue}\rule{\textwidth}{2.5pt}]
\titlespacing*{\chapter}{0pt}{-10pt}{30pt}

\titleformat{\section}
  {\normalfont\Large\bfseries\color{espritblue}}
  {{\color{espritgold}\thesection}}{0.6em}{}
  [\vspace{1pt}\color{espritblue!20}\rule{\textwidth}{0.5pt}]
\titlespacing*{\section}{0pt}{18pt}{6pt}

\titleformat{\subsection}
  {\normalfont\large\bfseries\color{espritblue!85}}
  {\thesubsection}{0.5em}{}
\titlespacing*{\subsection}{0pt}{12pt}{4pt}

\titleformat{\subsubsection}
  {\normalfont\normalsize\bfseries\color{espritblue!70}}
  {\thesubsubsection}{0.5em}{}
\titlespacing*{\subsubsection}{0pt}{8pt}{3pt}

%--- Caption style ------------------------------------------------------------
\captionsetup{font=small,labelfont={bf,color=espritblue},format=hang,margin=0.3cm}

%--- List spacing -------------------------------------------------------------
\setlist[itemize]{leftmargin=1.6em,itemsep=2pt,parsep=0pt,topsep=4pt}
\setlist[enumerate]{leftmargin=1.8em,itemsep=2pt,parsep=0pt,topsep=4pt}

%--- Table of contents styling ------------------------------------------------
\contentsmargin{0.5cm}
\titlecontents{chapter}[0em]
  {\addvspace{8pt}\large\bfseries\color{espritblue}}
  {\thecontentslabel\enspace}
  {}
  {\hfill\contentspage}
  [\vspace{2pt}\color{espritgold!70}\rule{\linewidth}{0.5pt}\vspace{2pt}]
\titlecontents{section}[1.8em]
  {\addvspace{2pt}\color{espritblue!80}}
  {\thecontentslabel\enspace}
  {}
  {\titlerule*[0.45pc]{.}\contentspage}
\titlecontents{subsection}[3.6em]
  {\addvspace{1pt}\small\color{espritblue!60}}
  {\thecontentslabel\enspace}
  {}
  {\titlerule*[0.45pc]{.}\contentspage}

%--- Useful macros -------------------------------------------------------------
\newcommand{\api}[1]{\texttt{#1}}
\newcommand{\code}[1]{\texttt{#1}}
\newcommand{\mqtt}[1]{\texttt{\textbf{#1}}}
\newcommand{\role}[1]{\textsl{#1}}
\newcommand{\done}{\ding{51}}
\newcommand{\notdone}{\ding{55}}
\newcommand{\sprint}[1]{\textbf{Sprint~#1}}

%=============================================================================
\begin{document}

%=============================================================================
%  TITLE PAGE
%=============================================================================
\begin{titlepage}
\thispagestyle{empty}
\begin{tikzpicture}[remember picture,overlay]
  % Top bar
  \fill[espritblue] (current page.north west) rectangle
    ([yshift=-3.5cm]current page.north east);
  % Gold accent line
  \fill[espritgold] ([yshift=-3.5cm]current page.north west) rectangle
    ([yshift=-3.8cm]current page.north east);
  % Bottom bar
  \fill[espritblue] (current page.south west) rectangle
    ([yshift=2cm]current page.south east);
  \fill[espritgold] ([yshift=2cm]current page.south west) rectangle
    ([yshift=2.3cm]current page.south east);
\end{tikzpicture}

\vspace*{0.5cm}

{\color{white}
\centering
{\large \textbf{ISET Siliana -- Institut Supérieur des Etudes Technologiques de Siliana}}\par
\vspace{0.2cm}
{\normalsize Département Informatique}\par
}

\vspace{3.5cm}

\centering
\begin{tcolorbox}[
  enhanced,
  colback=white,
  colframe=espritblue,
  arc=6pt,
  boxrule=2pt,
  left=1cm,right=1cm,top=0.8cm,bottom=0.8cm,
  drop shadow
]
\centering
{\color{espritblue}\Huge\textbf{Plateforme IoT Maison Intelligente}}\par
\vspace{0.4cm}
{\color{gray}\large Système Complet d'Automatisation Résidentielle Intelligente}\par
\vspace{0.6cm}
\hrule height 1pt
\vspace{0.6cm}
{\Large\textbf{Rapport de Projet de Fin d'Études (PFE)}}\par
\vspace{0.3cm}
{\normalsize Licence Appliquée en Informatique}
\end{tcolorbox}

\vspace{1.2cm}

{\color{espritgold}\rule{0.45\textwidth}{0.8pt}}
\hspace{0.5em}{\color{espritblue!40}\small\textbf{INFORMATIONS}}%
\hspace{0.5em}{\color{espritgold}\rule{0.45\textwidth}{0.8pt}}

\vspace{0.8cm}

\begin{tabular}{r@{\hspace{0.6em}:\hspace{0.8em}}l}
  \color{espritblue}\textbf{Réalisé par}
    & Imen Khadhrawi \\[3pt]
    & Malak Hammami \\[8pt]
  \color{espritblue}\textbf{Encadrant académique}
    & [Nom de l'Encadrant] \\[8pt]
  \color{espritblue}\textbf{Encadrant industriel}
    & [Nom du Tuteur Industriel] \\[8pt]
  \color{espritblue}\textbf{Année académique}
    & 2025 -- 2026 \\
\end{tabular}

\vspace{1.2cm}

{\color{espritgold}\rule{\textwidth}{0.8pt}}

\vspace{0.4cm}

{\small\color{gray}
\textit{Node.js \(\cdot\) React \(\cdot\) MongoDB \(\cdot\) MQTT \(\cdot\) ESP32 \(\cdot\) Socket.IO \(\cdot\) Groq AI}
}
\end{titlepage}

%=============================================================================
%  FRONT MATTER
%=============================================================================
\pagenumbering{roman}
\setcounter{page}{1}

\chapter*{Dédicace}
\addcontentsline{toc}{chapter}{Dédicace}
\begin{center}
\vspace{2cm}
{\large\itshape
À ma famille, dont le soutien inconditionnel a rendu ce voyage possible.\\[0.5cm]
À mes professeurs et mentors à l'ISET Siliana,\\
dont les conseils ont forgé mon état d'esprit d'ingénieur.\\[0.5cm]
À tous les développeurs qui croient que la technologie\\
peut rendre la vie quotidienne plus sûre et plus intelligente.
}
\end{center}

\chapter*{Remerciements}
\addcontentsline{toc}{chapter}{Remerciements}
Je tiens à exprimer ma profonde gratitude à mon encadrant académique pour ses conseils continus et ses retours perspicaces tout au long de ce projet. Je suis également reconnaissant envers mon tuteur industriel pour m'avoir fourni une perspective concrète sur les systèmes de maison intelligente et l'architecture IoT.

Je remercie chaleureusement le corps enseignant de l'ISET Siliana pour les solides fondations théoriques et pratiques fournies tout au long du cursus. Les cours en systèmes distribués, génie logiciel et systèmes embarqués ont été directement déterminants dans la conception et la mise en \oe{}uvre de cette plateforme.

Je remercie également mes collègues et amis qui ont participé aux sessions de test, signalé des bogues et suggéré des améliorations. Leurs retours pratiques ont considérablement amélioré l'expérience utilisateur du produit final.

Enfin, merci à la communauté open source derrière Node.js, React, MongoDB, Aedes MQTT et toutes les bibliothèques utilisées dans ce projet.

\chapter*{Résumé}
\addcontentsline{toc}{chapter}{Résumé}
Les systèmes de maison intelligente sont de plus en plus intégraux à la vie résidentielle moderne, mais la plupart des solutions commerciales restent fragmentées entre différentes applications, protocoles et fournisseurs. Ce rapport de PFE présente la \textbf{Plateforme IoT Maison Intelligente}, un système unifié, full-stack et de qualité production développé selon la méthodologie agile Scrum en six sprints.

La plateforme fournit un backend \textbf{Node.js/Express} avec un courtier \textbf{Aedes MQTT} intégré, une application web progressive \textbf{React/Vite}, et un firmware \textbf{ESP32} pour l'intégration des dispositifs physiques. Les fonctionnalités clés comprennent : la gestion multi-rôles (Agence, Administrateur, Résident), le contrôle des appareils en temps réel via MQTT et Socket.IO, un assistant IA alimenté par l'API Groq, la gestion des urgences gaz avec contrôle matériel (vanne et alarme), une visualisation 3D isométrique du plan d'étage, la surveillance énergétique, l'automatisation par scénarios temporels, et un système complet de journalisation des audits.

Le système a été conçu selon les principes d'architecture en couches, de modélisation orientée domaine et de communication événementielle. Les fonctionnalités de sécurité comprennent l'authentification JWT, Google OAuth 2.0, l'authentification à deux facteurs TOTP, la limitation de débit et les middlewares d'autorisation basés sur les rôles.

Ce rapport documente le cycle de vie Scrum complet : définition du backlog produit, planification des sprints, modélisation UML, implémentation, tests et déploiement, en concluant par une revue de sécurité et une feuille de route future.

\textbf{Mots-clés :} IoT, Maison Intelligente, MQTT, Node.js, React, MongoDB, ESP32, Socket.IO, Scrum, Agile, JWT, Groq IA, Systèmes Temps Réel.

\chapter*{Abstract}
\addcontentsline{toc}{chapter}{Abstract}
Smart home systems are increasingly integral to modern residential life, yet most commercial solutions remain fragmented. This PFE report presents the \textbf{SmartHome IoT Platform}, a unified, full-stack home automation system built using the Scrum agile methodology over six sprints. The platform integrates Node.js/Express, Aedes MQTT, React/Vite, Socket.IO, MongoDB, ESP32 firmware, and Groq AI, delivering multi-role governance, real-time device control, gas safety emergency handling, 3D floorplan visualization, and scenario automation.

\textbf{Keywords:} IoT, Smart Home, MQTT, Node.js, React, MongoDB, ESP32, Scrum, JWT, Groq AI.

\tableofcontents
\listoffigures
\listoftables

\clearpage
\pagenumbering{arabic}
\setcounter{page}{1}

%=============================================================================
%  CHAPTERS
%=============================================================================
%=============================================================================
%  Chapitre 1 — Introduction Générale
%=============================================================================
\chapter{Introduction Générale}

\section{Contexte et Motivation}

Les maisons connectées ont évolué d'écosystèmes de gadgets isolés vers des systèmes
cyber-physiques intégrés où la captation, l'actionnement, l'analyse et l'interaction
utilisateur sont étroitement couplés. Selon Statista (2024), le nombre de dispositifs
de maison intelligente en service dans le monde a dépassé 1,3 milliard d'unités, et le
marché mondial de la maison intelligente est projeté à 622 milliards de dollars d'ici 2026.

Malgré cette croissance, la plupart des solutions disponibles restent fragmentées : une
application mobile pour les caméras, une autre pour l'éclairage, et des outils distincts
pour les dispositifs de sécurité. Cette fragmentation crée des inefficacités opérationnelles,
affaiblit la posture de sécurité et génère des frictions pour les utilisateurs quotidiens.

La \textbf{Plateforme IoT Maison Intelligente} répond à cette lacune avec une solution
unifiée : une plateforme unique de qualité production qui intègre la gestion des appareils,
les canaux de commande et de télémétrie en temps réel, la gouvernance du foyer, les flux
d'urgence sécuritaire, un assistant IA, et des scénarios d'automatisation extensibles ---
le tout accessible depuis une interface web responsive.

\section{Cadre du Projet}

Ce projet a été développé en tant que Projet de Fin d'Études (PFE) à
\textbf{ISET Siliana --- Institut Supérieur des Etudes Technologiques de Siliana}, année
académique 2025--2026, en vue de l'obtention de la Licence Appliquée en Informatique.

Le projet est organisé comme un \textbf{monorepo} contenant trois artefacts déployables :

\begin{itemize}[leftmargin=1.5cm,itemsep=6pt]
  \item \textbf{Service backend} --- API REST Node.js/Express, courtier MQTT Aedes intégré,
        passerelle Socket.IO, et couche de persistance MongoDB.
  \item \textbf{Client frontend} --- Application monopage React~19/Vite avec routage
        conscient des rôles, synchronisation d'état en temps réel, plan d'étage 3D
        isométrique, et interface de chat IA.
  \item \textbf{Firmware ESP32} --- Sketch Arduino sur microcontrôleur ESP32-D intégrant
        la détection de température/humidité (DHT11), la détection de gaz (MQ6), le
        contrôle de LEDs RGB, et l'actionnement de servomoteurs.
\end{itemize}

\section{Problématique}

Le problème d'ingénierie central traité par ce projet est la conception et la mise en
œuvre d'une \textbf{plateforme multi-locataires de maison intelligente scalable} qui
satisfait simultanément les contraintes suivantes :

\begin{enumerate}[leftmargin=1.5cm,itemsep=6pt]
  \item \textbf{Réactivité temps réel} pour les événements utilisateur et la télémétrie
        IoT avec une latence bout en bout inférieure à 500~ms.
  \item \textbf{Séparation claire des rôles} entre les comptes Agence, Administrateur
        de maison, Résident et SuperAdmin avec une gestion granulaire des permissions.
  \item \textbf{Comportement critique de sécurité} pour les incidents gazeux avec des
        transitions d'état d'urgence déterministes, une actuation matérielle et une
        notification multi-canal.
  \item \textbf{Architecture maintenable} supportant l'extension progressive des
        fonctionnalités sans rompre les contrats existants.
  \item \textbf{Intégration IA accessible} permettant le contrôle des appareils en
        langage naturel via une API de modèle de langage à faible latence.
\end{enumerate}

\section{Objectifs du Projet}

Les objectifs du projet, mappés à des livrables mesurables, sont les suivants :

\begin{center}
\begin{longtable}{p{0.08\textwidth} p{0.48\textwidth} p{0.34\textwidth}}
\caption{Objectifs et Livrables du Projet}\label{tab:objectifs}\\
\toprule
\textbf{ID} & \textbf{Objectif} & \textbf{Livrable} \\
\midrule
\endfirsthead
\toprule
\textbf{ID} & \textbf{Objectif} & \textbf{Livrable} \\
\midrule
\endhead
\bottomrule
\endfoot
O1  & Construire un système d'authentification multi-rôles complet et sécurisé & API Auth, middleware JWT, OAuth, 2FA \\
O2  & Activer les flux CRUD et de commande pour les appareils connectés & API Appareils, pipeline de publication MQTT \\
O3  & Intégrer un courtier MQTT pour la messagerie IoT & Courtier Aedes, routage de topics, client ESP32 \\
O4  & Livrer la synchronisation d'état en temps réel aux clients web & Passerelle Socket.IO, catalogue d'événements \\
O5  & Implémenter la surveillance d'urgence gaz avec réponse matérielle & Service Gas Monitor, workflow d'urgence \\
O6  & Fournir un flux de demande de permission pour les résidents & API Permissions, système de notification \\
O7  & Construire un plan d'étage 3D avec contrôle interactif des appareils & Composant FloorPlan, overlay SVG \\
O8  & Développer un assistant IA (Melo) avec contrôle en langage naturel & API Companion, interface chat Melo \\
O9  & Implémenter l'automatisation par scénarios temporels & API Scénarios, planificateur node-cron \\
O10 & Créer un tableau de bord de surveillance de la consommation énergétique & API Énergie, visualisation recharts \\
O11 & Déployer le système complet avec runbook opérationnel documenté & Config déploiement, annexe runbook \\
\end{longtable}
\end{center}

\section{Approche Méthodologique --- Scrum}

La mise en œuvre a suivi le cadre agile \textbf{Scrum}, choisi pour ses cycles de livraison
itératifs, l'intégration continue du feedback et la gestion transparente du backlog.
Les rôles Scrum ont été mappés au contexte académique :

\begin{itemize}[leftmargin=1.5cm,itemsep=6pt]
  \item \textbf{Product Owner} --- Encadrant académique (validation des exigences fonctionnelles).
  \item \textbf{Scrum Master} --- Auteur (facilitation des cérémonies, suivi de la vélocité).
  \item \textbf{Équipe de développement} --- Auteur (développement full-stack et firmware).
\end{itemize}

\noindent Le projet a été livré en \textbf{six sprints} de deux semaines chacun, comme
résumé dans le Tableau~\ref{tab:sprint_resume}.

\begin{center}
\begin{longtable}{c p{0.26\textwidth} p{0.38\textwidth} c}
\caption{Résumé des Sprints}\label{tab:sprint_resume}\\
\toprule
\textbf{Sprint} & \textbf{Thème} & \textbf{Livrables Principaux} & \textbf{SP} \\
\midrule
\endfirsthead
\toprule
\textbf{Sprint} & \textbf{Thème} & \textbf{Livrables Principaux} & \textbf{SP} \\
\midrule
\endhead
\bottomrule
\endfoot
1 & Authentification et Gestion des Utilisateurs & Flux d'inscription, JWT, OAuth, 2FA, vérification email & 34 \\
2 & Gestion des Appareils et MQTT              & CRUD Appareils, courtier MQTT, firmware ESP32           & 39 \\
3 & Plan d'Étage et Synchronisation TR         & Plan 3D, scènes/ambiances, passerelle Socket.IO         & 37 \\
4 & Systèmes de Sécurité                       & Moniteur gaz, DHT11, workflow d'urgence, vanne          & 36 \\
5 & Assistant IA et Automatisation             & Melo IA, scénarios, surveillance énergétique            & 34 \\
6 & Portail Agence et Audit                    & Hub conciergerie, journaux d'audit, permissions         & 37 \\
\midrule
  & \textbf{Total}                             &                                                          & \textbf{217} \\
\bottomrule
\end{longtable}
\end{center}

\section{Contributions Scientifiques et Techniques}

Ce projet contribue au domaine de l'IoT et du génie logiciel web en plusieurs dimensions :

\begin{itemize}[leftmargin=1.5cm,itemsep=6pt]
  \item Une \textbf{architecture pratique} combinant APIs REST, courtage MQTT et WebSockets
        dans un service Node.js cohérent.
  \item Un \textbf{modèle de gouvernance des rôles} adaptable aux contextes résidentiels
        et de gestion de propriétés multi-unités.
  \item Une \textbf{boucle de contrôle sécuritaire} liant les seuils de télémétrie PPM
        aux commandes matérielles et aux notifications multi-canal.
  \item Une \textbf{interface en langage naturel} alimentée par LLM, contrainte par le
        contrôle d'accès basé sur les rôles.
  \item Un \textbf{plan d'artefacts Scrum} applicable aux projets IoT similaires en
        contexte académique.
\end{itemize}

\section{Organisation du Rapport}

Ce rapport est structuré pour guider le lecteur des aspects conceptuels aux détails
d'implémentation et opérationnels :

\begin{description}[leftmargin=1.5cm,itemsep=6pt]
  \item[Chapitres 1--3 :] Contexte du projet, analyse des besoins, état de l'art et
        évaluation technologique.
  \item[Chapitre 4 :] Architecture système et modélisation UML (déploiement, composants,
        cas d'utilisation, classes et diagrammes de séquence).
  \item[Chapitres 5--8 :] Implémentation sprint par sprint couvrant la base de données,
        les APIs backend, l'interface frontend et l'intégration IoT.
  \item[Chapitres 9--11 :] Évaluation sécuritaire, stratégie de test, pipeline de
        déploiement et gestion de projet.
  \item[Chapitre 12 :] Conclusion, rétrospective et feuille de route future.
  \item[Annexes (13--16) :] Référence API complète, inventaire frontend, runbook
        d'installation et matrice de traçabilité.
\end{description}

\chapter{Requirement Analysis and Specifications}

\section{Stakeholders and Roles}
The platform serves multiple stakeholders whose responsibilities define the authorization model:

\subsection{Agency}
Agency users supervise multiple housing units, review owner requests, monitor cross-unit status, and maintain governance controls at the platform level.

\subsection{House Owner / Admin}
House admins manage one household domain: they register devices, supervise residents, review permission requests, activate scenes, and operate safety controls.

\subsection{Resident}
Residents are authenticated users linked to a house code. Their capabilities are intentionally constrained and can be expanded via explicit admin approval.

\subsection{System Operator}
A technical operator (DevOps/maintainer) maintains infrastructure, observability, backups, and updates.

\section{Functional Requirements}

\subsection{Identity and Access}
\begin{itemize}[leftmargin=1.2cm]
  \item Registration and login for resident, admin, and agency flows.
  \item JWT-based API authentication.
  \item Google OAuth login integration.
  \item Email verification and password reset lifecycle.
  \item Route-level and action-level role authorization.
\end{itemize}

\subsection{House and Unit Lifecycle}
\begin{itemize}[leftmargin=1.2cm]
  \item Generation and validation of unit/house codes.
  \item Join existing managed unit or create solo unmanaged unit.
  \item Track claimed/unclaimed unit status.
\end{itemize}

\subsection{Device Management}
\begin{itemize}[leftmargin=1.2cm]
  \item Register IoT devices with metadata (topic, type, room, position).
  \item List and delete devices per house context.
  \item Send command payloads to devices via MQTT topics.
  \item Update online/offline state and last-seen timestamps.
\end{itemize}

\subsection{Permission Workflow}
\begin{itemize}[leftmargin=1.2cm]
  \item Resident submits restricted action request.
  \item Admin reviews and approves/denies.
  \item Approved permissions are merged into resident capabilities.
  \item Notification state supports unread/read transitions.
\end{itemize}

\subsection{Scenarios and Automation}
\begin{itemize}[leftmargin=1.2cm]
  \item Activation of named scenarios (Morning, Dinner, Cinematic, Sleep).
  \item Bulk update of lights and openings according to scene presets.
  \item Synchronization of dashboard state after scenario changes.
\end{itemize}

\subsection{Gas Safety}
\begin{itemize}[leftmargin=1.2cm]
  \item Ingest gas readings from MQTT sensor topics.
  \item Compare readings to configurable threshold.
  \item Trigger emergency mode, alarm command, and valve closure.
  \item Broadcast emergency state to relevant clients in real time.
  \item Allow admin clearing and threshold tuning.
\end{itemize}

\subsection{Realtime Communication}
\begin{itemize}[leftmargin=1.2cm]
  \item Socket.IO authenticated channels for user notifications.
  \item House-room broadcasts for dashboard synchronization.
  \item Agency-global alerts for emergency escalation.
\end{itemize}

\section{Non-Functional Requirements}

\subsection{Performance}
The system should provide low-latency updates for state-changing events, with response times suitable for human-perceived immediacy in dashboard interactions.

\subsection{Scalability}
Architecture should support horizontal evolution in API and socket layers, and future broker scaling strategies.

\subsection{Security}
Critical requirements include password hashing, signed tokens, role isolation, transport security readiness, and controlled data exposure.

\subsection{Maintainability}
Code must be modular, readable, and extensible. Route handlers, models, middlewares, and real-time components should follow clear boundaries.

\subsection{Reliability}
Failure in one communication channel should not corrupt persistent state; safety-critical transitions should remain deterministic and recoverable.

\section{Use Case Inventory}
\begin{longtable}{p{3cm}p{4cm}p{7cm}}
\toprule
Actor & Use Case & Description \\
\midrule
Agency & Review owner requests & Approve/reject account creation requests and issue one-time access code. \\
Agency & Monitor units & View claimed and unclaimed units, statuses, and analytics indicators. \\
Admin & Register device & Add a new smart device and map it to room/topic. \\
Admin & Activate scenario & Apply global scene parameters to eligible devices. \\
Admin & Review request & Approve or deny resident permission request. \\
Admin & Clear emergency & Manually clear emergency mode after hazard assessment. \\
Resident & Request permission & Ask access to restricted actions/rooms. \\
Resident & Control assigned room & Execute allowed actions within authorization scope. \\
System & Process gas reading & Persist reading, evaluate threshold, emit events, publish hardware commands. \\
\bottomrule
\end{longtable}

\section{Acceptance Criteria}
Each major capability has measurable criteria:
\begin{itemize}[leftmargin=1.2cm]
  \item Authentication endpoints return valid token and role payload on success.
  \item Device commands are published to expected MQTT topic format.
  \item Permission approval updates resident permissions and emits real-time update.
  \item Gas reading above threshold transitions to emergency and emits alert events.
  \item Role guards prevent unauthorized access to protected endpoints/pages.
\end{itemize}

\section{Constraints and Assumptions}
\begin{itemize}[leftmargin=1.2cm]
  \item Initial development targets local network and localhost orchestration.
  \item MQTT payload schemas are controlled by platform-owned device firmware simulations.
  \item Users have internet/email access for verification workflows.
  \item Hardware-level commands are simulated in development and can map to physical devices in deployment.
\end{itemize}

\section{Requirement Traceability Overview}
A traceability matrix is maintained in Appendix chapters, mapping each requirement to API routes, model entities, and frontend pages. This ensures auditable engineering alignment from specification to implementation.

\chapter{State of the Art and Technology Stack}

\section{Smart Home Platforms: Current Landscape}
Smart-home systems generally follow one of three patterns:
\begin{enumerate}[leftmargin=1.2cm]
  \item Vendor-locked ecosystems focused on proprietary devices.
  \item Gateway-centric open-source stacks with limited user governance.
  \item Cloud-native IoT services with strong scaling but high integration overhead.
\end{enumerate}

The SmartHome platform targets a hybrid engineering path: flexible protocol-level integration (MQTT), application-level governance (REST + role model), and web-native UX (React).

\section{Protocol and Interaction Model}
\subsection{MQTT for Device Messaging}
MQTT is lightweight, topic-based, and suitable for constrained devices and low-bandwidth communication. The project uses MQTT for sensor telemetry and device command delivery.

\subsection{HTTP/REST for Business Operations}
REST endpoints handle account management, CRUD workflows, permissions, and administrative functions where request-response semantics are natural.

\subsection{Socket.IO for User-facing Real-time Events}
Socket.IO complements REST by delivering asynchronous updates such as approval notifications, gas alerts, and dashboard state changes.

\section{Technology Selection Rationale}
\subsection{Backend: Node.js + Express}
Node.js provides asynchronous I/O efficiency for mixed workloads (HTTP, sockets, MQTT events). Express offers low ceremony and direct control over middleware and route composition.

\subsection{Database: MongoDB + Mongoose}
MongoDB supports document flexibility for evolving IoT payload fields and role-specific user metadata. Mongoose adds schema constraints and lifecycle hooks (e.g., password hashing).

\subsection{MQTT Broker: Aedes}
Aedes allows embedding broker behavior directly inside backend runtime, reducing integration complexity in prototyping and early production stages.

\subsection{Frontend: React + Vite}
React enables component-level decomposition and route guards. Vite ensures fast development iteration and efficient modern build tooling.

\subsection{Realtime: Socket.IO}
Socket.IO handles transport fallback and namespace/room abstractions, simplifying per-house and per-role event fan-out.

\section{Comparative Positioning}
\begin{table}[H]
\centering
\caption{Comparative perspective of architectural choices}
\begin{tabular}{p{4cm}p{4cm}p{6cm}}
\toprule
Dimension & Typical Alternative & Adopted Choice and Advantage \\
\midrule
Device Messaging & HTTP polling from devices & MQTT topics: lower overhead, event-native pattern \\
Realtime UI & Frequent client polling & Socket push model: lower latency and less load \\
Data Model & Strict relational schema & Flexible documents for IoT and user metadata \\
Auth Model & Single-role simple auth & Multi-role with permission workflow and onboarding \\
\bottomrule
\end{tabular}
\end{table}

\section{Engineering Standards and Principles}
The project follows practical engineering principles:
\begin{itemize}[leftmargin=1.2cm]
  \item Separation of concerns across route, model, middleware, and realtime services.
  \item Explicit role checks for privileged operations.
  \item Deterministic state transitions for safety-critical workflows.
  \item Event-driven synchronization between data persistence and UI state.
\end{itemize}

\section{Toolchain and Dependencies}
\subsection{Backend Core Libraries}
Important dependencies include Express, Mongoose, Aedes, Socket.IO server, JWT, bcrypt, and OAuth strategy packages.

\subsection{Frontend Core Libraries}
Important dependencies include React, React Router, Axios, Socket.IO client, MQTT client, and chart/animation utilities.

\subsection{Development Tooling}
Nodemon and Vite support development cycles; ESLint supports code quality baselines.

\section{Known Trade-offs}
\begin{itemize}[leftmargin=1.2cm]
  \item Embedded broker simplifies architecture but couples API and broker runtime.
  \item Flexible schemas ease evolution but require strict validation discipline.
  \item Socket-heavy UX improves responsiveness but demands robust auth handling and room hygiene.
\end{itemize}

\section{Chapter Synthesis}
The selected stack is coherent with project constraints: quick iteration, protocol-level interoperability, and real-time operational responsiveness. Subsequent chapters translate these choices into formal architecture and implementation.

\chapter{System Architecture and UML Modeling}

\section{Architectural Vision}
The SmartHome platform follows a layered event-driven architecture combining:
\begin{itemize}[leftmargin=1.2cm]
  \item API-oriented business services (REST).
  \item Event-oriented machine communication (MQTT).
  \item Real-time user synchronization (Socket.IO).
  \item Persistent state and history (MongoDB).
\end{itemize}

This architecture separates concerns while allowing fast propagation of state changes across subsystems.

\section{High-Level Context Diagram}
\begin{figure}[H]
\centering
\resizebox{\textwidth}{!}{%
\begin{tikzpicture}[
  node distance=1.2cm and 1.6cm,
  actor/.style={draw, rounded corners, fill=blue!10, minimum width=3.2cm, minimum height=0.9cm, align=center},
  system/.style={draw, rounded corners, fill=green!12, minimum width=5.0cm, minimum height=1.2cm, align=center},
  datastore/.style={draw, cylinder, shape border rotate=90, aspect=0.2, minimum height=1.2cm, minimum width=2.0cm, fill=orange!15, align=center}
]
\node[actor] (agency) {Agency User};
\node[actor, below=of agency] (admin) {House Admin};
\node[actor, below=of admin] (resident) {Resident};

\node[system, right=3.2cm of admin] (platform) {SmartHome Platform\\(API + Realtime + Broker)};
\node[datastore, right=3.6cm of platform] (db) {MongoDB};
\node[actor, below=1.6cm of platform] (devices) {IoT Devices};

\draw[-{Latex[length=3mm]}] (agency) -- node[above,sloped]{Manage units, requests} (platform);
\draw[-{Latex[length=3mm]}] (admin) -- node[above,sloped]{Control home, review permissions} (platform);
\draw[-{Latex[length=3mm]}] (resident) -- node[below,sloped]{Use allowed features} (platform);
\draw[-{Latex[length=3mm]}] (platform) -- node[above]{Persist/read} (db);
\draw[-{Latex[length=3mm]}] (devices) -- node[right]{MQTT telemetry / commands} (platform);
\end{tikzpicture}
}
\caption{System context diagram}
\end{figure}

\section{Component Diagram}
\begin{figure}[H]
\centering
\resizebox{\textwidth}{!}{%
\begin{tikzpicture}[
  comp/.style={draw, rounded corners, fill=cyan!10, minimum width=3.4cm, minimum height=1.0cm, align=center},
  arr/.style={-{Latex[length=3mm]}}
]
\node[comp] (frontend) at (0,0) {React Frontend};
\node[comp] (api) at (5,0) {Express API Layer};
\node[comp] (socket) at (10,0) {Socket.IO Service};

\node[comp] (auth) at (5,-2) {Auth and RBAC Middleware};
\node[comp] (domain) at (5,-4) {Domain Routes and Services};
\node[comp] (broker) at (10,-2) {Aedes MQTT Broker};
\node[comp] (gas) at (10,-4) {Gas Monitor Service};
\node[comp] (mongo) at (5,-6) {MongoDB and Models};

\draw[arr] (frontend) -- node[above]{REST} (api);
\draw[arr] (frontend) -- node[above]{Realtime events} (socket);
\draw[arr] (api) -- (auth);
\draw[arr] (auth) -- (domain);
\draw[arr] (domain) -- (mongo);
\draw[arr] (broker) -- node[right]{publish/subscribe} (gas);
\draw[arr] (gas) -- (mongo);
\draw[arr] (gas) -- node[left]{emit alerts} (socket);
\draw[arr] (domain) -- (broker);
\draw[arr] (api) -- (socket);
\end{tikzpicture}
}
\caption{Logical component architecture}
\end{figure}

\section{Deployment Diagram}
\begin{figure}[H]
\centering
\resizebox{\textwidth}{!}{%
\begin{tikzpicture}[
  box/.style={draw, rounded corners, fill=gray!10, minimum width=4.4cm, minimum height=1.2cm, align=center},
  arr/.style={-{Latex[length=3mm]}}
]
\node[box] (browser) at (0,0) {Client Browser\\(React SPA)};
\node[box] (node) at (6,0) {Node.js Runtime\\Express + Socket.IO + Aedes};
\node[box] (mongo) at (12,0) {MongoDB Server};
\node[box] (iot) at (6,-2.2) {IoT Devices / Simulators\\MQTT TCP / MQTT over WS};
\node[box] (mail) at (12,-2.2) {Email Service\\(verification/reset)};

\draw[arr] (browser) -- node[above]{HTTPS/HTTP + WS} (node);
\draw[arr] (node) -- node[above]{Mongoose} (mongo);
\draw[arr] (iot) -- node[left]{MQTT topics} (node);
\draw[arr] (node) -- node[right]{SMTP/API} (mail);
\end{tikzpicture}
}
\caption{Deployment-level topology}
\end{figure}

\section{Use Case Diagram Narrative}
The main use-case clusters are:
\begin{itemize}[leftmargin=1.2cm]
  \item Agency governance: manage units, review owner requests, monitor analytics.
  \item Admin operations: manage users and devices, approve permissions, activate scenes.
  \item Resident interactions: controlled dashboard usage, permission requests, profile actions.
  \item Automated safety loop: ingest sensor values, compute emergency state, propagate alerts.
\end{itemize}

\section{Sequence Diagram: Resident Permission Request}
\begin{figure}[H]
\centering
\resizebox{\textwidth}{!}{%
\begin{tikzpicture}[
  lifeline/.style={draw, minimum width=2.6cm, minimum height=0.8cm, fill=blue!8, align=center},
  arr/.style={-{Latex[length=2.5mm]}}
]
\node[lifeline] (res) at (0,0) {Resident UI};
\node[lifeline] (api) at (4,0) {Permissions API};
\node[lifeline] (db) at (8,0) {MongoDB};
\node[lifeline] (sock) at (12,0) {Socket.IO};
\node[lifeline] (adm) at (16,0) {Admin UI};

\draw[dashed] (res) -- (0,-5.2);
\draw[dashed] (api) -- (4,-5.2);
\draw[dashed] (db) -- (8,-5.2);
\draw[dashed] (sock) -- (12,-5.2);
\draw[dashed] (adm) -- (16,-5.2);

\draw[arr] (0,-1) -- node[above]{POST request} (4,-1);
\draw[arr] (4,-2) -- node[above]{create doc} (8,-2);
\draw[arr] (8,-2.6) -- node[below]{created} (4,-2.6);
\draw[arr] (4,-3.4) -- node[above]{emit event} (12,-3.4);
\draw[arr] (12,-4.2) -- node[above]{requestCreated} (16,-4.2);
\end{tikzpicture}
}
\caption{Permission request event propagation}
\end{figure}

\section{Sequence Diagram: Gas Emergency}
\begin{figure}[H]
\centering
\resizebox{\textwidth}{!}{%
\begin{tikzpicture}[
  lifeline/.style={draw, minimum width=2.4cm, minimum height=0.8cm, fill=red!8, align=center},
  arr/.style={-{Latex[length=2.5mm]}}
]
\node[lifeline] (sensor) at (0,0) {Gas Sensor};
\node[lifeline] (broker) at (3.8,0) {MQTT Broker};
\node[lifeline] (monitor) at (7.6,0) {Gas Monitor};
\node[lifeline] (db) at (11.4,0) {HouseState DB};
\node[lifeline] (sock) at (15.2,0) {Socket.IO};
\node[lifeline] (ui) at (19,0) {Dashboards};

\foreach \x in {0,3.8,7.6,11.4,15.2,19} {\draw[dashed] (\x,-0.4) -- (\x,-6.2);} 

\draw[arr] (0,-1.0) -- node[above]{publish sensors/H1234/gas} (3.8,-1.0);
\draw[arr] (3.8,-1.8) -- node[above]{event packet} (7.6,-1.8);
\draw[arr] (7.6,-2.6) -- node[above]{save level + emergencyMode} (11.4,-2.6);
\draw[arr] (7.6,-3.4) -- node[above]{publish alarm + valve command} (3.8,-3.4);
\draw[arr] (7.6,-4.2) -- node[above]{emit house:gas-emergency} (15.2,-4.2);
\draw[arr] (15.2,-5.0) -- node[above]{real-time alert} (19,-5.0);
\end{tikzpicture}
}
\caption{Gas emergency sequence}
\end{figure}

\section{Activity Diagram Narrative: User Onboarding}
User onboarding branch logic:
\begin{enumerate}[leftmargin=1.2cm]
  \item User authenticates.
  \item If house code already exists in profile, redirect to dashboard.
  \item Else choose between join existing unit and create solo unit.
  \item Persist onboarding result and assign role capabilities.
  \item Navigate to role-appropriate dashboard.
\end{enumerate}

\section{State Diagram Narrative: House Safety State}
House safety states can be conceptualized as:
\begin{itemize}[leftmargin=1.2cm]
  \item Normal: gas level below threshold, emergency disabled.
  \item Alert Triggered: gas level exceeds threshold and transition detected.
  \item Emergency Active: emergency mode true, valve closed, alarm command retained.
  \item Recovery: value returns below threshold or admin clear action executed.
\end{itemize}

\section{Architectural Quality Attributes}
The architecture optimizes for:
\begin{itemize}[leftmargin=1.2cm]
  \item Modularity: independent route and service modules.
  \item Responsiveness: push events to clients instead of polling.
  \item Evolvability: model and route extensions with localized impact.
  \item Operational clarity: deterministic handling for safety events.
\end{itemize}

\section{Design Constraints and Evolution Strategy}
The architecture currently co-hosts API, sockets, and broker in one process. This is acceptable for project and early deployment stages. A future production strategy can separate broker, API, and realtime gateway into independent services while preserving contracts defined in this chapter.

\chapter{Database Design and Domain Models}

\section{Data Strategy}
MongoDB is used as the primary persistence layer. The design follows bounded domain entities with pragmatic denormalization where needed for query simplicity and runtime performance.

\section{Main Collections}

\subsection{User}
The user model captures identity and role governance:
\begin{itemize}[leftmargin=1.2cm]
  \item Core identity: name, email, password hash.
  \item Role domain: admin, resident, agency.
  \item House linkage: houseCode, linkedAdmin.
  \item Access metadata: permissions, assignedRoom.
  \item Account lifecycle: email verification and reset token fields.
  \item OAuth metadata: googleId.
\end{itemize}

Unique constraints and generation logic enforce house-code integrity for admin accounts.

\subsection{Device}
Devices represent connected assets and include:
\begin{itemize}[leftmargin=1.2cm]
  \item Human metadata: name, room, type, color.
  \item Integration metadata: topic, owner, houseCode.
  \item Runtime state: status, state payload, brightness/openPct.
  \item Spatial metadata: position coordinates for floor plan rendering.
\end{itemize}

\subsection{HouseState}
House-level global state includes:
\begin{itemize}[leftmargin=1.2cm]
  \item Lighting mode and occupancy/security flags.
  \item Gas safety data: gasLevel, gasThreshold, emergencyMode, gasValveOpen.
  \item Audit fields: updatedBy and timestamps.
\end{itemize}

\subsection{PermissionRequest}
This collection implements access-governance workflow:
\begin{itemize}[leftmargin=1.2cm]
  \item requester and houseAdmin references.
  \item action key/label and optional room scope.
  \item status lifecycle: pending, approved, denied.
  \item read-state metadata for admin and requester.
\end{itemize}

\subsection{Unit}
Unit entities support agency-level management:
\begin{itemize}[leftmargin=1.2cm]
  \item unique unitCode, claim status, and owner linkage.
  \item managed/unmanaged mode with operational status indicators.
\end{itemize}

\subsection{HouseOwnerRequest}
This collection tracks candidate owner requests submitted for agency review, including approval workflow and one-time access code lifecycle.

\section{Entity Relationship View}
\begin{figure}[H]
\centering
\begin{tikzpicture}[
  entity/.style={draw, rounded corners, fill=purple!8, minimum width=3.2cm, minimum height=1cm, align=center},
  rel/.style={-{Latex[length=3mm]}}
]
\node[entity] (user) at (0,0) {User};
\node[entity] (device) at (5,0) {Device};
\node[entity] (house) at (10,0) {HouseState};
\node[entity] (perm) at (0,-3) {PermissionRequest};
\node[entity] (unit) at (5,-3) {Unit};
\node[entity] (ownerreq) at (10,-3) {HouseOwnerRequest};

\draw[rel] (user) -- node[above]{owns/manages} (device);
\draw[rel] (device) -- node[above]{belongs to} (house);
\draw[rel] (user) -- node[left]{creates/reviews} (perm);
\draw[rel] (user) -- node[above,sloped]{linked by houseCode} (house);
\draw[rel] (unit) -- node[right]{owned by} (user);
\draw[rel] (ownerreq) -- node[right]{approved to create} (user);
\end{tikzpicture}
\caption{Conceptual data relationship view}
\end{figure}

\section{Schema Decisions and Justification}
\subsection{Document-oriented Flexibility}
Device state payloads may differ by type; MongoDB mixed fields simplify this heterogeneity.

\subsection{House Code as Domain Anchor}
Many queries are scoped by houseCode, making it a key partition-like attribute in authorization and event fan-out.

\subsection{Auditability}
Timestamp and reviewer fields across workflow collections provide traceability for governance actions.

\section{Indexing Strategy}
Recommended indexes for scale:
\begin{itemize}[leftmargin=1.2cm]
  \item User(email) unique.
  \item User(role, houseCode) for role/house queries.
  \item Device(topic) unique and Device(houseCode, room).
  \item PermissionRequest(houseCode, status, createdAt).
  \item Unit(unitCode) unique.
\end{itemize}

\section{Data Integrity Rules}
\begin{itemize}[leftmargin=1.2cm]
  \item Password must be hashed before persistence.
  \item House code pattern must follow one-letter + four-digit format.
  \item Permission request decisions restricted to approved/denied set.
  \item Gas threshold bounded to safe operational range.
\end{itemize}

\section{Migration and Evolution Notes}
Future schema evolution can include:
\begin{itemize}[leftmargin=1.2cm]
  \item Time-series collection for telemetry historical analytics.
  \item Event journal for full audit replay.
  \item Soft-delete flags for compliance and operational recovery.
\end{itemize}

\section{Data Governance Summary}
The current model set balances expressive domain coverage and implementation simplicity while preserving clear migration paths for advanced analytics and enterprise compliance.

\chapter{Backend Design and API Specification}

\section{Backend Runtime Overview}
The backend process orchestrates three communication surfaces in one Node.js runtime:
\begin{itemize}[leftmargin=1.2cm]
  \item Express HTTP API for transactional and business workflows.
  \item Socket.IO server for authenticated push notifications.
  \item Aedes MQTT broker for machine telemetry and command channels.
\end{itemize}

\section{Bootstrapping Pipeline}
System startup sequence:
\begin{enumerate}[leftmargin=1.2cm]
  \item Load environment variables.
  \item Initialize Express middlewares (JSON parser, CORS, passport).
  \item Register route modules.
  \item Create HTTP server and bind Socket.IO.
  \item Bind MQTT-over-WebSocket endpoint.
  \item Connect MongoDB.
  \item Start gas monitor listener.
  \item Listen on HTTP and MQTT TCP ports.
\end{enumerate}

\section{Middleware Layer}
\subsection{Authentication}
JWT middleware extracts Bearer token, verifies signature, loads user context, and attaches it to request object.

\subsection{Authorization}
Role gates enforce admin-only and agency-only operations, reducing privilege escalation risk.

\subsection{Input Validation}
Core route handlers validate required fields, normalized codes, and decision enums. Future enhancement should centralize validation schemas.

\section{Route Domains}

\subsection{Authentication and Account Lifecycle}
Includes resident/admin/agency login/register, email verification, password reset, OAuth callback, and legacy compatibility handlers.

\subsection{Unit and Agency Management}
Supports unit generation, claim/join workflow, owner oversight, and aggregated agency analytics.

\subsection{Device Management}
Admin-managed registration, listing scoped by house filter, command emission via MQTT topics, and deletion lifecycle.

\subsection{Scenario Operations}
Scenario activation endpoint updates house state and device states in bulk and then emits synchronized dashboard update.

\subsection{Permission Governance}
Resident request submission and admin decision endpoints with notification lifecycle and permission-merge behavior.

\subsection{Gas Safety Endpoints}
State retrieval, manual threshold updates, valve control, and emergency clearing for admin users.

\section{API Contract Design Principles}
\begin{itemize}[leftmargin=1.2cm]
  \item House-scoped filtering by authenticated user context.
  \item Deterministic JSON payload structure.
  \item Meaningful status codes for business and authorization outcomes.
  \item Route naming consistent with domain concepts.
\end{itemize}

\section{Representative Endpoint Matrix}
\begin{longtable}{p{2.3cm}p{4.2cm}p{3.2cm}p{5.2cm}}
\toprule
Method & Endpoint & Access & Purpose \\
\midrule
POST & /api/auth/resident/register & Public & Create resident account and issue token \\
POST & /api/auth/admin/login & Public & Authenticate admin with access code \\
POST & /api/units/join & Authenticated & Link user to existing unit code \\
GET & /api/devices & Authenticated & Retrieve house-scoped devices \\
POST & /api/devices & Admin & Register new device \\
PUT & /api/devices/:id/command & Authenticated & Publish command payload to device topic \\
POST & /api/permissions/request & Resident & Request restricted action approval \\
PATCH & /api/permissions/admin/requests/:id & Admin & Approve or deny request \\
GET & /api/gas/state & Authenticated & Read current gas status \\
PATCH & /api/gas/threshold & Admin & Update emergency threshold \\
\bottomrule
\end{longtable}

\section{Error Handling Strategy}
Error handling currently uses route-level try/catch patterns with JSON error messages. For production hardening, a centralized error middleware is recommended to enforce a uniform error schema and correlation IDs.

\section{Backend Security Controls}
\begin{itemize}[leftmargin=1.2cm]
  \item Password hashing via bcrypt.
  \item Signed JWT for stateless auth.
  \item Role-based endpoint restrictions.
  \item Token-verified socket handshake.
  \item Controlled account lifecycle with email verification and password reset.
\end{itemize}

\section{Observability Hooks}
The current implementation provides practical console-level observability for connection events, MQTT publishes, and critical transitions. Recommended next step is structured logging and metrics export.

\section{Backend Evolution Opportunities}
\begin{itemize}[leftmargin=1.2cm]
  \item Introduce schema-based validation layer.
  \item Add rate limiting and request throttling.
  \item Split broker and API services for independent scaling.
  \item Add integration test harness for route-event coupling.
\end{itemize}

%=============================================================================
%  Chapitre 7 — Sprint 4 : Systèmes de Sécurité Gaz et DHT11
%=============================================================================
\chapter{Sprint 4 --- Détection Gaz, Capteur DHT11 et Protocoles d'Urgence}

\begin{tcolorbox}[sprintbox,title={Sprint 4 --- Systèmes de Sécurité}]
\textbf{Durée :} 2 semaines \quad
\textbf{Points de Story :} 36 \quad
\textbf{Objectif :} Implémenter la détection d'urgence gaz de bout en bout depuis le
capteur MQ6 jusqu'à l'actuation matérielle et l'alerte UI plein écran, plus la
surveillance de température et humidité DHT11.
\end{tcolorbox}

\section{Service de Fond Gas Monitor}

Le service \texttt{gasMonitor.js} est le cerveau sécuritaire de la plateforme. Il
s'abonne à tous les topics de capteurs gaz, évalue les lectures PPM par rapport au seuil
configuré, et pilote la machine à états d'urgence.

\begin{lstlisting}[style=codeblock,language=JavaScript,caption={Service Gas Monitor (backend/gasMonitor.js)}]
const initGasMonitor = (mqttClient, io) => {
  mqttClient.subscribe('sensors/+/gas');

  mqttClient.on('message', async (topic, message) => {
    const parts = topic.split('/');
    if (parts.length !== 3 || parts[2] !== 'gas') return;
    const houseCode = parts[1];

    let ppm;
    try {
      ppm = JSON.parse(message.toString()).ppm;
    } catch { return; }   // Ignorer les messages malformes

    let state = await HouseState.findOne({ houseCode });
    if (!state) state = await HouseState.create({ houseCode });

    const threshold = state.gas.threshold ?? 400;
    state.gas.level = ppm;
    await state.save();

    if (ppm >= threshold && !state.emergencyMode) {
      // ---- ACTIVATION DE L'URGENCE ----
      state.emergencyMode  = true;
      state.gas.lastAlert  = new Date();
      state.gas.valveOpen  = false;
      await state.save();

      // 1. Alerter tous les clients web de la maison
      io.to(houseCode).emit('house:gas-emergency', { ppm, threshold, houseCode });

      // 2. Fermer la vanne gaz
      mqttClient.publish(`command/${houseCode}/valve`, 'CLOSE', { qos: 1 });

      // 3. Activer le buzzer
      mqttClient.publish(`command/${houseCode}/alarm`, 'ON', { qos: 1 });

      // 4. Journaliser l'audit
      await AuditLog.create({
        actorName: 'Systeme',
        actorRole: 'system',
        action: `Urgence gaz : ${ppm.toFixed(0)} PPM depasse le seuil ${threshold}`,
        category: 'security',
        houseCode
      });
    }
  });
};
\end{lstlisting}

\section{Calcul PPM du Capteur Gaz MQ6}

Le capteur MQ6 produit une tension analogique proportionnelle à la concentration de gaz.
L'ESP32 lit cette valeur via l'ADC et la convertit en PPM en utilisant une formule
calibrée :

\begin{lstlisting}[style=codeblock,language=C++,caption={Calcul PPM Capteur Gaz (ESP32)}]
float readGasPPM() {
  int raw      = analogRead(GAS_PIN);           // 0 - 4095 (ADC 12 bits)
  float voltage= raw * 3.3f / 4095.f;
  float Rs     = ((3.3f - voltage) / voltage) * 10.0f;  // RL = 10 kOhm
  float ratio  = Rs / 9.83f;                            // calibration Ro
  return 1000.0f * pow(ratio, -2.37f);                  // PPM approximation
}

// Publier toutes les 3 secondes
if (millis() - lastGasPublish >= 3000) {
  float ppm = readGasPPM();
  String payload = "{\"ppm\":" + String(ppm, 1) + "}";
  mqttClient.publish(gasTopic.c_str(), payload.c_str());

  // Securite locale : declencher le buzzer si ADC brut > 800
  if (analogRead(GAS_PIN) > 800) tone(BUZZER_PIN, 1000);
  else                           noTone(BUZZER_PIN);
  lastGasPublish = millis();
}
\end{lstlisting}

\section{Service DHT11 et Firmware}

\begin{lstlisting}[style=codeblock,language=C++,caption={Lecture DHT11 et Publication MQTT (ESP32)}]
#include <DHT.h>
DHT dht(DHT_PIN, DHT11);

// Dans loop() :
if (millis() - lastDHTPublish >= 10000) {
  float temperature = dht.readTemperature();
  float humidity    = dht.readHumidity();

  if (!isnan(temperature) && !isnan(humidity)) {
    String payload = "{\"temperature\":" + String(temperature, 1)
                   + ",\"humidity\":"    + String(humidity,    1) + "}";
    mqttClient.publish(dhtTopic.c_str(), payload.c_str());
  }
  lastDHTPublish = millis();
}
\end{lstlisting}

\section{Overlay d'Urgence Frontend}

\begin{lstlisting}[style=codeblock,language=JavaScript,caption={Composant GasEmergencyOverlay}]
const GasEmergencyOverlay = ({ ppm, onClear, isAdmin }) => (
  <motion.div
    className="fixed inset-0 z-[9999] bg-red-600/90 flex flex-col
               items-center justify-center text-white"
    animate={{ opacity: [1, 0.8, 1] }}
    transition={{ repeat: Infinity, duration: 1.2 }}
  >
    <AlertTriangle size={96} className="mb-6 animate-bounce" />
    <h1 className="text-5xl font-bold mb-4">FUITE DE GAZ DETECTEE</h1>
    <p className="text-2xl mb-2">Niveau actuel : {ppm?.toFixed(0)} PPM</p>
    <p className="text-lg mb-8 opacity-80">
      Vanne fermee - Alarme active - Evacuez immediatement
    </p>
    {isAdmin && (
      <button onClick={onClear}
        className="px-8 py-3 bg-white text-red-700 rounded-full
                   font-bold text-lg hover:bg-red-100 transition">
        Lever l'Urgence
      </button>
    )}
  </motion.div>
);
\end{lstlisting}

\section{Revue et Rétrospective Sprint 4}

\begin{tcolorbox}[goalbox,title={Revue Sprint 4}]
\textbf{Complété :} 6 éléments sur 6 (36 PS). Urgence gaz vérifiée sur le matériel. \\
\textbf{Points forts :} Souffler sur le capteur MQ6 a déclenché l'overlay d'urgence en
500~ms, activé le buzzer physique, et le servo de vanne s'est fermé. L'admin a levé
l'urgence depuis le tableau de bord et le buzzer s'est tu en 1 seconde. \\
\textbf{Rétrospective :}
\begin{itemize}[leftmargin=1cm,itemsep=2pt]
  \item \textbf{Ce qui a bien fonctionné :} L'architecture de sécurité à deux couches
        (déclencheur local ESP32 ADC~$>$~800 + seuil serveur 400~PPM) fournit une
        défense en profondeur.
  \item \textbf{Défi :} Le capteur DHT11 retourne occasionnellement NaN lors des lectures
        à froid ; géré par un guard \texttt{isnan()} avant la publication MQTT.
  \item \textbf{Amélioration :} Ajout de \texttt{qos:1} aux commandes alarme et vanne
        pour survivre aux déconnexions MQTT brèves lors des scénarios d'urgence.
\end{itemize}
\end{tcolorbox}

%=============================================================================
%  Chapitre 8 — Sprint 5 : Compagnon IA Melo, Scénarios et Surveillance Énergétique
%=============================================================================
\chapter{Sprint 5 --- Compagnon IA (Melo), Automatisation par Scénarios et Surveillance Énergétique}

\begin{tcolorbox}[sprintbox,title={Sprint 5 --- Compagnon IA \& Automatisation}]
\textbf{Durée :} 2 semaines \quad
\textbf{Points de Story :} 34 \quad
\textbf{Objectif :} Livrer le compagnon IA Melo pour le contrôle des appareils en langage
naturel, l'automatisation temporelle via node-cron, et un tableau de bord de surveillance
de la consommation énergétique.
\end{tcolorbox}

\section{Backlog Sprint 5}

\begin{center}
\begin{longtable}{p{0.07\textwidth} p{0.54\textwidth} p{0.08\textwidth} p{0.12\textwidth}}
\caption{Décomposition des User Stories Sprint 5}\label{tab:sprint5_backlog}\\
\toprule
\textbf{ID} & \textbf{User Story} & \textbf{PS} & \textbf{Statut} \\
\midrule
\endfirsthead
\toprule
\textbf{ID} & \textbf{User Story} & \textbf{PS} & \textbf{Statut} \\
\midrule
\endhead
\bottomrule
\endfoot
US-23 & Compagnon IA (Melo) pour contrôle d'appareils en langage naturel & 13 & \done Terminé \\
US-24 & Scénarios d'automatisation temporels avec node-cron & 8 & \done Terminé \\
US-25 & Graphiques de consommation énergétique horaire & 8 & \done Terminé \\
US-26 & Historique et statut des demandes de permission & 5 & \done Terminé \\
\end{longtable}
\end{center}

\section{Compagnon IA Melo}

\subsection{Architecture}

Melo est construit sur l'\textbf{API Groq}, un service d'inférence cloud offrant des temps
de réponse LLM inférieurs à la seconde. Le backend expose un seul endpoint
\api{POST /api/companion/chat} qui :

\begin{enumerate}[leftmargin=1.5cm,itemsep=4pt]
  \item Enrichit le message utilisateur avec un prompt système structuré contenant la liste
        des appareils courants, les lectures des capteurs et les permissions utilisateur.
  \item Envoie l'historique de conversation à l'API Groq avec le modèle
        \texttt{llama-3.3-70b-versatile}.
  \item Analyse la réponse LLM pour détecter les \textbf{blocs d'action} (commandes JSON structurées).
  \item Exécute les actions validées via le même pipeline de commande d'appareils que l'API REST.
  \item Retourne la réponse LLM et les résultats des actions au frontend.
\end{enumerate}

\subsection{Construction du Prompt Système}

\begin{lstlisting}[style=codeblock,language=JavaScript,caption={Construction du Prompt Syst\`{e}me Melo (backend/routes/companion.js)}]
const buildSystemPrompt = (user, devices, houseState) => {
  const accessibleDevices = user.role === 'admin'
    ? devices
    : devices.filter(d =>
        d.room === user.room ||
        user.permissions.includes(`room:${d.room.toLowerCase().replace(' ', '-')}`) ||
        user.permissions.includes('global:controls')
      );

  return `
Tu es Melo, un assistant maison intelligente intelligent et amical.
Tu aides les utilisateurs a controler leurs appareils en langage naturel.

UTILISATEUR COURANT :
- Nom : ${user.name}
- Role : ${user.role}
- Piece assignee : ${user.room || 'Toutes les pieces (admin)'}

APPAREILS ACCESSIBLES (${accessibleDevices.length} au total) :
${accessibleDevices.map(d =>
  `- ${d.name} (${d.type}) dans ${d.room} | topic: ${d.mqttTopic} | etat: ${d.state}`
).join('\n')}

ENVIRONNEMENT :
- Niveau gaz : ${houseState?.gas?.level?.toFixed(0) ?? 'N/A'} PPM
- Temperature : ${houseState?.temperature ?? 'N/A'}C
- Humidite : ${houseState?.humidity ?? 'N/A'}%
- Scene courante : ${houseState?.lightingMode ?? 'Aucune'}

CAPACITES :
- Pour controler un appareil, reponds avec un bloc ACTION :
  [ACTION: {"deviceId":"<id>","command":{"state":"ON"|"OFF","brightness":0-100}}]
- Pour activer une scene : [SCENE: "Morning"|"Dinner"|"Cinematic"|"Sleep"]
- Pour des requetes d'information, reponds de maniere conversationnelle.
- Ne controle JAMAIS des appareils hors de la liste accessible.
- Si tu ne peux pas faire quelque chose (permission refusee), explique poliment.
`.trim();
};
\end{lstlisting}

\subsection{Implémentation de l'Endpoint Chat}

\begin{lstlisting}[style=codeblock,language=JavaScript,caption={Endpoint Chat Compagnon (backend/routes/companion.js)}]
router.post('/chat', protect, async (req, res) => {
  const { messages } = req.body;   // Historique de conversation du frontend

  const devices    = await Device.find({ houseCode: req.user.houseCode });
  const houseState = await HouseState.findOne({ houseCode: req.user.houseCode });
  const sysPrompt  = buildSystemPrompt(req.user, devices, houseState);

  // Appel API Groq
  const completion = await groq.chat.completions.create({
    model: 'llama-3.3-70b-versatile',
    messages: [
      { role: 'system', content: sysPrompt },
      ...messages
    ],
    temperature: 0.3,              // Basse temperature pour des commandes deterministes
    max_tokens: 1024
  });

  const reply = completion.choices[0].message.content;

  // Analyser et executer les blocs ACTION
  const actionMatches = [...reply.matchAll(/\[ACTION: ({.*?})\]/g)];
  const results = [];
  for (const match of actionMatches) {
    try {
      const { deviceId, command } = JSON.parse(match[1]);
      const device = devices.find(d => d._id.toString() === deviceId);
      if (device) {
        await executeDeviceCommand(device, command, mqttClient, io, req.user);
        results.push({ deviceId, success: true });
      }
    } catch { /* Ignorer les blocs d'action malformes */ }
  }

  // Analyser les blocs SCENE
  const sceneMatch = reply.match(/\[SCENE: "([^"]+)"\]/);
  if (sceneMatch) {
    await activateScene(sceneMatch[1], req.user.houseCode, mqttClient, io);
  }

  res.json({ reply, actionsExecuted: results.length });
});
\end{lstlisting}

\subsection{Composant Frontend Melo}

L'interface Melo est un widget de chat flottant rendu dans le coin inférieur droit de
l'application. Elle inclut une mascotte 3D animée (\texttt{Melo3D}) construite avec
React Three Fiber et animée avec Framer Motion.

\begin{lstlisting}[style=codeblock,language=JavaScript,caption={Interface Chat AuraCompanion (frontend)}]
const AuraCompanion = () => {
  const [open, setOpen]       = useState(false);
  const [messages, setMessages]= useState([]);
  const [input, setInput]     = useState('');
  const [loading, setLoading] = useState(false);
  const { token }             = useAuth();

  const sendMessage = async () => {
    const userMsg = { role: 'user', content: input };
    setMessages(prev => [...prev, userMsg]);
    setInput('');
    setLoading(true);
    try {
      const { data } = await axios.post('/api/companion/chat',
        { messages: [...messages, userMsg] },
        { headers: { Authorization: `Bearer ${token}` } });
      setMessages(prev => [
        ...prev, { role: 'assistant', content: data.reply }
      ]);
    } finally {
      setLoading(false);
    }
  };

  return (
    <div className="fixed bottom-6 right-6 z-50">
      <AnimatePresence>
        {open && (
          <motion.div
            initial={{ scale: 0, opacity: 0 }}
            animate={{ scale: 1, opacity: 1 }}
            exit={{ scale: 0, opacity: 0 }}
            className="w-80 h-[450px] bg-white dark:bg-gray-900
                       rounded-2xl shadow-2xl flex flex-col overflow-hidden"
          >
            <div className="bg-indigo-600 p-3 text-white font-bold
                            flex items-center gap-2">
              <Sparkles size={18} /> Melo -- Votre IA Maison Intelligente
            </div>
            <div className="flex-1 overflow-y-auto p-3 space-y-3">
              {messages.map((m, i) => (
                <div key={i} className={`flex
                  ${m.role === 'user' ? 'justify-end' : 'justify-start'}`}>
                  <div className={`rounded-xl px-3 py-2 max-w-[80%] text-sm
                    ${m.role === 'user'
                      ? 'bg-indigo-100'
                      : 'bg-gray-100 dark:bg-gray-800'}`}>
                    {m.content
                      .replace(/\[ACTION:.*?\]/g, '')
                      .replace(/\[SCENE:.*?\]/g, '')}
                  </div>
                </div>
              ))}
              {loading && (
                <div className="text-xs text-gray-400 animate-pulse">
                  Melo reflechit...
                </div>
              )}
            </div>
            <div className="p-3 border-t flex gap-2">
              <input value={input} onChange={e => setInput(e.target.value)}
                onKeyDown={e => e.key === 'Enter' && sendMessage()}
                className="flex-1 rounded-lg border px-3 py-1.5 text-sm"
                placeholder="Posez une question a Melo..." />
              <button onClick={sendMessage}
                className="bg-indigo-600 text-white rounded-lg px-3 py-1.5 text-sm">
                Envoyer
              </button>
            </div>
          </motion.div>
        )}
      </AnimatePresence>
      <button onClick={() => setOpen(!open)}
        className="w-14 h-14 rounded-full bg-indigo-600 text-white shadow-lg
                   flex items-center justify-center hover:bg-indigo-700 transition">
        <Bot size={28} />
      </button>
    </div>
  );
};
\end{lstlisting}

\section{Modèle de Scénario et Moteur d'Automatisation}

\subsection{Schéma Scénario}

\begin{lstlisting}[style=codeblock,language=JavaScript,caption={Sch\'ema Sc\'enario (backend/models/Scenario.js)}]
const scenarioSchema = new mongoose.Schema({
  name:         { type: String, required: true },
  description:  { type: String },
  houseCode:    { type: String, required: true },
  isActive:     { type: Boolean, default: true },
  triggerType:  { type: String, enum: ['time'], default: 'time' },
  scheduleTime: { type: String, required: true },  // Format "HH:MM"
  action:       { type: String, required: true },  // ex. "mood:Morning"
  targetDevice: { type: mongoose.Schema.Types.ObjectId, ref: 'Device' },
  lastTriggeredAt: { type: Date }
}, { timestamps: true });
\end{lstlisting}

\subsection{Service Planificateur}

\begin{lstlisting}[style=codeblock,language=JavaScript,caption={Planificateur de Sc\'enarios (backend/scheduler.js)}]
const { schedule } = require('node-cron');

const initScheduler = (mqttClient, io) => {
  // Executer chaque minute
  schedule('* * * * *', async () => {
    const now = new Date();
    const currentTime = `${String(now.getHours()).padStart(2,'0')}:`
                      + `${String(now.getMinutes()).padStart(2,'0')}`;

    const scenarios = await Scenario.find({
      isActive: true,
      scheduleTime: currentTime
    });

    for (const scenario of scenarios) {
      try {
        if (scenario.action.startsWith('mood:')) {
          const scene = scenario.action.split(':')[1];
          await activateScene(scene, scenario.houseCode, mqttClient, io);
        } else if (scenario.action === 'lock_all') {
          await lockAllDoors(scenario.houseCode, mqttClient);
        } else if (scenario.action.startsWith('toggle:light:')) {
          const state = scenario.action.endsWith(':on') ? 'ON' : 'OFF';
          await toggleLights(scenario.houseCode, state, mqttClient);
        }
        scenario.lastTriggeredAt = new Date();
        await scenario.save();

        io.to(scenario.houseCode).emit('scenario:triggered', {
          name: scenario.name, action: scenario.action
        });

        await AuditLog.create({
          actorName: 'Planificateur',
          actorRole: 'system',
          action: `Scenario "${scenario.name}" declenche : ${scenario.action}`,
          category: 'scenario',
          houseCode: scenario.houseCode
        });
      } catch (err) {
        console.error(`[Planificateur] Erreur scenario ${scenario._id}:`, err);
      }
    }
  });
};
\end{lstlisting}

\section{Surveillance Énergétique}

\subsection{Algorithme de Calcul de la Consommation}

La consommation énergétique est estimée à partir de la collection AuditLog. Chaque
événement ON pour un appareil marque le début d'une période de consommation. L'événement
OFF correspondant (ou la fin de session) la ferme. Le kWh pour chaque période est :
\[ \text{kWh} = \frac{W_{\text{appareil}} \times \Delta t_{\text{heures}}}{1000} \]

\begin{lstlisting}[style=codeblock,language=JavaScript,caption={Endpoint \'Energie (backend/routes/energy.js)}]
router.get('/hourly', protect, async (req, res) => {
  const WATTAGE = { light:10, lamp:10, 'rgb-light':10, ac:1500,
                    tv:80, fan:50, outlet:200, 'washing-machine':500 };

  const logs = await AuditLog.find({
    houseCode: req.user.houseCode,
    category:  'device',
    createdAt: { $gte: new Date(Date.now() - 86400000) }  // Dernieres 24h
  }).sort({ createdAt: 1 });

  // Construire la carte de consommation horaire
  const hourly = Array(24).fill(0);
  const openPeriods = {};  // deviceId -> { startHour, watt }

  for (const log of logs) {
    const deviceId = log.metadata?.deviceId;
    const action   = log.action;
    const hour     = new Date(log.createdAt).getHours();
    const watt     = WATTAGE[log.metadata?.deviceType] ?? 10;

    if (action.includes('ON') || action.includes('OPEN')) {
      openPeriods[deviceId] = { startHour: hour, watt };
    } else if (openPeriods[deviceId]) {
      const duration = (hour - openPeriods[deviceId].startHour) || 1;
      hourly[hour] += (openPeriods[deviceId].watt * duration) / 1000;
      delete openPeriods[deviceId];
    }
  }

  res.json({ hourly, totalKwh: hourly.reduce((a,b) => a+b, 0) });
});
\end{lstlisting}

\section{Système de Demandes de Permission}

\subsection{Modèle DemandePermission}

\begin{lstlisting}[style=codeblock,language=JavaScript,caption={Sch\'ema DemandePermission (backend/models/PermissionRequest.js)}]
const permissionRequestSchema = new mongoose.Schema({
  requester:   { type: mongoose.Schema.Types.ObjectId, ref: 'User', required: true },
  houseAdmin:  { type: mongoose.Schema.Types.ObjectId, ref: 'User', required: true },
  actionKey:   { type: String, required: true },   // ex. 'room:kitchen'
  actionLabel: { type: String, required: true },   // libelle lisible
  room:        { type: String },                   // piece demandee (optionnel)
  status:      { type: String,
                 enum: ['pending','approved','denied'], default: 'pending' },
  reviewer:    { type: mongoose.Schema.Types.ObjectId, ref: 'User' },
  reviewNotes: { type: String },
  adminReadAt:     { type: Date },
  requesterReadAt: { type: Date }
}, { timestamps: true });
\end{lstlisting}

\subsection{Flux d'Approbation de Permission}

\begin{lstlisting}[style=codeblock,language=JavaScript,caption={Approbation Demande de Permission (backend/routes/permissions.js)}]
router.post('/:id/approve', protect, admin, async (req, res) => {
  const { reviewNotes } = req.body;
  const request = await PermissionRequest.findById(req.params.id)
    .populate('requester');

  if (!request || request.status !== 'pending')
    return res.status(400).json({ message: 'Demande introuvable ou deja traitee' });

  request.status      = 'approved';
  request.reviewer    = req.user._id;
  request.reviewNotes = reviewNotes;
  await request.save();

  // Accorder la permission au resident
  const resident = await User.findById(request.requester._id);
  if (!resident.permissions.includes(request.actionKey)) {
    resident.permissions.push(request.actionKey);
  }
  if (request.room && !resident.room) {
    resident.room = request.room;
  }
  await resident.save();

  // Notifier le resident en temps reel
  io.to(`user:${resident._id}`).emit('permission:updated', {
    status: 'approved', actionKey: request.actionKey
  });

  res.json({ success: true });
});
\end{lstlisting}

\section{Revue et Rétrospective Sprint 5}

\begin{tcolorbox}[goalbox,title={Revue Sprint 5}]
\textbf{Complété :} 4 stories sur 4 (34 PS). La démo Melo a impressionné. \\
\textbf{Points forts :} Taper ``Allume toutes les lumières de la chambre à 80\%'' a fait
illuminer la LED RGB chambre à 80\% de luminosité en 1,5 secondes. Le graphique d'énergie
a montré des pics de consommation réalistes lors du déclenchement du scénario climatisation. \\
\textbf{Rétrospective :}
\begin{itemize}[leftmargin=1cm,itemsep=2pt]
  \item \textbf{Ce qui a bien fonctionné :} La latence d'inférence Groq ($<$800~ms) rend
        Melo très réactif.
  \item \textbf{Défi :} L'analyse des blocs ACTION a nécessité des regex soigneuses pour
        éviter les faux positifs dans les réponses conversationnelles.
  \item \textbf{Amélioration :} Ajout de \texttt{temperature: 0.3} pour réduire les noms
        d'appareils hallucités dans les réponses LLM.
\end{itemize}
\end{tcolorbox}

%=============================================================================
%  Chapitre 9 — Sprint 6 : Portail Agence, Journal d'Audit et Finalisation
%=============================================================================
\chapter{Sprint 6 --- Portail Agence, Journal d'Audit et Finalisation}

\begin{tcolorbox}[sprintbox,title={Sprint 6 --- Portail Agence et Audit}]
\textbf{Durée :} 2 semaines \quad
\textbf{Points de Story :} 37 \quad
\textbf{Objectif :} Livrer le hub conciergerie agence avec streaming en direct des maisons,
un tableau de bord de journal d'audit complet, les demandes de permission résidents,
les modes de sécurité et la gestion du profil.
\end{tcolorbox}

\section{Hub Conciergerie et Portail Agence}

\subsection{Server-Sent Events (SSE) pour le Statut en Direct}

\begin{lstlisting}[style=codeblock,language=JavaScript,caption={Endpoint SSE pour le Hub Conciergerie (backend/routes/concierge.js)}]
router.get('/stream', protect, agencyOnly, (req, res) => {
  res.setHeader('Content-Type',  'text/event-stream');
  res.setHeader('Cache-Control', 'no-cache');
  res.setHeader('Connection',    'keep-alive');
  res.flushHeaders();

  const push = (data) => res.write(`data: ${JSON.stringify(data)}\n\n`);

  // Envoyer l'instantane initial de tous les statuts admin
  push({ type: 'snapshot', houses: getOnlineAdmins() });

  // S'abonner aux evenements de connexion admin
  const onConnect    = (admin) => push({ type:'online',  houseCode: admin.houseCode });
  const onDisconnect = (admin) => push({ type:'offline', houseCode: admin.houseCode });

  adminEvents.on('connect',    onConnect);
  adminEvents.on('disconnect', onDisconnect);

  // Nettoyage a la deconnexion du client
  req.on('close', () => {
    adminEvents.off('connect',    onConnect);
    adminEvents.off('disconnect', onDisconnect);
    res.end();
  });
});
\end{lstlisting}

\section{Système de Journal d'Audit}

\subsection{Schéma AuditLog}

\begin{lstlisting}[style=codeblock,language=JavaScript,caption={Sch\'ema AuditLog (backend/models/AuditLog.js)}]
const auditLogSchema = new mongoose.Schema({
  actor:      { type: mongoose.Schema.Types.ObjectId, ref: 'User' },
  actorName:  { type: String, required: true },
  actorRole:  { type: String },
  action:     { type: String, required: true },
  category:   { type: String,
                enum: ['auth','device','security','user','scenario','system'],
                default: 'system' },
  details:    { type: String },
  metadata:   { type: mongoose.Schema.Types.Mixed },
  ipAddress:  { type: String },
  houseCode:  { type: String }
}, { timestamps: true });
\end{lstlisting}

\subsection{Événements Journalisés}

\begin{center}
\begin{longtable}{p{0.16\textwidth} p{0.72\textwidth}}
\caption{Catégories et Événements du Journal d'Audit}\label{tab:audit_events}\\
\toprule
\textbf{Catégorie} & \textbf{Événements Journalisés} \\
\midrule
\endfirsthead
\toprule
\textbf{Catégorie} & \textbf{Événements Journalisés} \\
\midrule
\endhead
\bottomrule
\endfoot
\texttt{auth}     & Connexion (succès/échec), inscription, réinitialisation mdp, vérification email, 2FA activer/désactiver, connexion Google OAuth \\
\texttt{device}   & Appareil créé, mis à jour, supprimé, commande exécutée (ON/OFF/luminosité/scène) \\
\texttt{security} & Urgence gaz activée, urgence levée, mode Absence activé, mode Verrouillage activé \\
\texttt{user}     & Résident ajouté/supprimé, rôle modifié, permission accordée/révoquée \\
\texttt{scenario} & Scénario créé/mis à jour/supprimé, scénario déclenché par le planificateur \\
\texttt{system}   & Démarrage serveur, démarrage courtier MQTT, initialisation planificateur \\
\end{longtable}
\end{center}

\section{Modes de Sécurité}

\begin{center}
\begin{longtable}{p{0.18\textwidth} p{0.70\textwidth}}
\caption{Définition des Modes de Sécurité}\label{tab:security_modes}\\
\toprule
\textbf{Mode} & \textbf{Comportement} \\
\midrule
\endfirsthead
\toprule
\textbf{Mode} & \textbf{Comportement} \\
\midrule
\endhead
\bottomrule
\endfoot
Normal           & Fonctionnement standard ; tous les appareils librement contrôlables par les utilisateurs autorisés. \\
Mode Absence     & Surveillance caméra renforcée ; toutes les portes extérieures verrouillées ; fenêtres fermées ; vanne gaz fermée ; événements de mouvement journalisés. \\
Mode Verrouillage & Sécurité maximale ; toutes les caméras armées ; toutes les portes et fenêtres verrouillées ; vanne gaz fermée ; contrôle des appareils résidents suspendu. \\
\end{longtable}
\end{center}

\section{Récapitulatif Complet des Routes API}

\begin{center}
\begin{longtable}{p{0.26\textwidth} p{0.14\textwidth} p{0.48\textwidth}}
\caption{Toutes les Routes Backend API}\label{tab:api_summary}\\
\toprule
\textbf{Préfixe de Route} & \textbf{Auth} & \textbf{Rôle} \\
\midrule
\endfirsthead
\toprule
\textbf{Préfixe de Route} & \textbf{Auth} & \textbf{Rôle} \\
\midrule
\endhead
\bottomrule
\endfoot
\api{/api/auth}         & Aucune/JWT   & Inscription, connexion, vérification email, réinitialisation mdp \\
\api{/api/oauth}        & Aucune       & Callback Google OAuth 2.0 \\
\api{/api/2fa}          & JWT          & Configuration, vérification, désactivation TOTP 2FA \\
\api{/api/profile}      & JWT          & Obtenir/mettre à jour le profil et l'avatar utilisateur \\
\api{/api/users}        & Admin JWT    & Liste, mise à jour, suppression des utilisateurs de la maison \\
\api{/api/devices}      & Admin JWT    & CRUD des appareils \\
\api{/api/floorplan}    & JWT          & État plan, commandes, scènes \\
\api{/api/scenarios}    & Admin JWT    & CRUD des scénarios \\
\api{/api/gas}          & Admin JWT    & État gaz, seuil, levée urgence \\
\api{/api/permissions}  & JWT          & CRUD demandes de permission et approbation \\
\api{/api/messages}     & JWT          & Soumission de message support et réponse admin \\
\api{/api/agency}       & Agence JWT   & Liste des admins pour la vue agence \\
\api{/api/units}        & Admin JWT    & Gestion des unités/maisons \\
\api{/api/concierge}    & Agence JWT   & Données hub conciergerie et flux SSE \\
\api{/api/audit}        & Admin JWT    & Requête et pagination du journal d'audit \\
\api{/api/stats}        & Admin JWT    & Statistiques tableau de bord \\
\api{/api/energy}       & Admin JWT    & Consommation énergétique horaire/quotidienne \\
\api{/api/companion}    & JWT          & Endpoint chat assistant IA \\
\api{/api/owner-requests}& Aucune/JWT  & Demandes de compte propriétaire \\
\end{longtable}
\end{center}

\section{Revue et Rétrospective Sprint 6}

\begin{tcolorbox}[goalbox,title={Revue Sprint 6}]
\textbf{Complété :} 8 stories sur 8 (37 PS). Plateforme fonctionnellement complète. \\
\textbf{Points forts :} Le hub conciergerie agence a affiché 3 maisons simulées avec le
clignotement en ligne/hors ligne en direct. Le journal d'audit a été exporté avec succès
en PDF avec jsPDF. Le mode Absence a immédiatement verrouillé tous les servomoteurs de
porte sur l'ESP32. \\
\textbf{Rétrospective :}
\begin{itemize}[leftmargin=1cm,itemsep=2pt]
  \item \textbf{Ce qui a bien fonctionné :} SSE était plus simple que Socket.IO pour le
        cas d'utilisation de streaming en lecture seule de l'agence.
  \item \textbf{Défi :} Le formatage de l'export PDF a nécessité un réglage minutieux des
        largeurs de colonnes dans jsPDF pour les longues chaînes d'action d'audit.
  \item \textbf{Amélioration :} Ajout d'une limite de 2~Mo sur les uploads d'avatar après
        qu'une image base64 de 15~Mo a causé des avertissements de taille de document MongoDB.
\end{itemize}
\end{tcolorbox}

\chapter{Deployment, DevOps, and Operations}

\section{Environment Topology}
A practical environment model includes:
\begin{itemize}[leftmargin=1.2cm]
  \item Development: local monorepo with backend and frontend concurrently.
  \item Staging: cloud VM or container stack with test data.
  \item Production: hardened setup with managed database and observability stack.
\end{itemize}

\section{Configuration Management}
Environment variables control ports, database URI, JWT secret, OAuth credentials, frontend origin, and access codes. Secret values must never be committed to source control.

\section{Build and Run Pipeline}
\subsection{Root Orchestration}
The monorepo root starts backend and frontend together for integrated development testing.

\subsection{Backend Runtime}
Backend service starts HTTP server, websocket endpoint, MQTT broker, and DB connection in one process.

\subsection{Frontend Runtime}
Vite serves SPA with hot-module reload in development and optimized static build in production.

\section{Containerization Strategy (Recommended)}
A production-ready containerization path:
\begin{enumerate}[leftmargin=1.2cm]
  \item Backend image with Node runtime, non-root user, healthcheck endpoint.
  \item Frontend static build served by Nginx or edge hosting.
  \item Externalized MongoDB as managed database service.
\end{enumerate}

\section{CI/CD Blueprint}
\begin{itemize}[leftmargin=1.2cm]
  \item Trigger on pull request and main branch merge.
  \item Install dependencies and run lint checks.
  \item Run automated test suites.
  \item Build artifacts and publish deploy package.
  \item Deploy to staging, then production with approval gate.
\end{itemize}

\section{Observability Architecture}
\subsection{Logs}
Adopt structured logs with request IDs and event correlation IDs.

\subsection{Metrics}
Capture HTTP latency, broker throughput, socket connection counts, and emergency incidents.

\subsection{Alerts}
Configure alerting for service downtime, DB connectivity issues, and abnormal emergency burst patterns.

\section{Backup and Recovery}
\begin{itemize}[leftmargin=1.2cm]
  \item Daily MongoDB backups with retention policy.
  \item Configuration backup for environment templates.
  \item Recovery drill procedure tested periodically.
\end{itemize}

\section{Operational Runbook Essentials}
Operational teams should maintain runbooks for:
\begin{itemize}[leftmargin=1.2cm]
  \item Incident triage and escalation.
  \item Broker connection troubleshooting.
  \item Emergency false alert investigation.
  \item User account recovery and access revocation.
\end{itemize}

\section{Scalability Roadmap}
\begin{itemize}[leftmargin=1.2cm]
  \item Separate broker and API services.
  \item Introduce Redis adapter for horizontal socket scaling.
  \item Add telemetry queueing for burst management.
  \item Partition data and house-level workloads.
\end{itemize}

\section{Operational KPIs}
\begin{table}[H]
\centering
\caption{Sample operational service-level indicators}
\begin{tabular}{p{5cm}p{3cm}p{6cm}}
\toprule
Indicator & Target & Notes \\
\midrule
API Availability & > 99.5\% & Monthly rolling target \\
p95 API Latency & < 350 ms & Authenticated business routes \\
Socket Delivery Success & > 99\% & Event emission to intended rooms \\
Emergency Processing Delay & < 2 s & Sensor receipt to user alert \\
\bottomrule
\end{tabular}
\end{table}

\section{Operations Summary}
The project is deployment-ready for educational and pilot usage. With incremental hardening in CI/CD, observability, and scaling controls, it can evolve into production-grade service operations.

\chapter{Project Management, Planning, and Risk Governance}

\section{Project Planning Approach}
The project followed iterative delivery with milestone-based checkpoints:
\begin{itemize}[leftmargin=1.2cm]
  \item Requirement analysis and actor mapping.
  \item Architectural baseline and domain modeling.
  \item Backend feature slices with route and model implementation.
  \item Frontend role flows and dashboard composition.
  \item Realtime integration and safety workflows.
  \item Validation, documentation, and reporting.
\end{itemize}

\section{Work Breakdown Structure}
\begin{longtable}{p{2cm}p{5cm}p{7cm}}
\toprule
Phase & Work Package & Deliverables \\
\midrule
P1 & Analysis and specification & Requirements document, use cases, acceptance criteria \\
P2 & Architecture and UML & Context/component/deployment diagrams, sequence models \\
P3 & Backend core implementation & Auth, users, units, devices, permissions, gas APIs \\
P4 & Frontend implementation & Routing guards, dashboards, contexts, pages \\
P5 & Realtime and IoT integration & Socket room model, MQTT broker flow, emergency loop \\
P6 & Validation and deployment prep & Tests, hardening checklist, deployment notes \\
P7 & Report finalization & Full PFE manuscript and appendices \\
\bottomrule
\end{longtable}

\section{Timeline (Illustrative)}
\begin{table}[H]
\centering
\caption{Milestone timeline}
\begin{tabular}{p{5cm}p{4cm}p{5cm}}
\toprule
Milestone & Period & Outcome \\
\midrule
M1: Scope freeze & Week 1-2 & Functional perimeter approved \\
M2: Architecture freeze & Week 3-4 & Core technical decisions validated \\
M3: Backend MVP & Week 5-8 & Essential APIs and models delivered \\
M4: Frontend MVP & Week 9-11 & Role-based UI and route guards delivered \\
M5: Realtime and safety & Week 12-14 & MQTT + Socket workflows operational \\
M6: Validation and report & Week 15-18 & Testing, optimization, and full report \\
\bottomrule
\end{tabular}
\end{table}

\section{Risk Register}
\begin{longtable}{p{4cm}p{2cm}p{2cm}p{6cm}}
\toprule
Risk & Probability & Impact & Mitigation \\
\midrule
Missing critical dependency in runtime & Medium & High & Lock dependency list, perform clean install checks \\
Auth misconfiguration between frontend and backend & Medium & High & Contract tests and environment templates \\
MQTT payload inconsistency & Medium & Medium & Introduce schema validation and topic contract docs \\
Role leakage in endpoints & Low & High & Strict middleware and security test cases \\
Realtime event drift & Medium & Medium & Event naming registry and client compatibility tests \\
\bottomrule
\end{longtable}

\section{Communication and Governance}
Weekly status reviews and integration demos reduce delivery risk by exposing mismatches early between UI assumptions and backend behavior.

\section{Cost and Resource Perspective}
The chosen open-source stack minimizes licensing cost and is suitable for academic projects and small teams. Main investment remains engineering time and operational discipline.

\section{Lessons Learned}
\begin{itemize}[leftmargin=1.2cm]
  \item Event contracts must be documented as first-class artifacts.
  \item Role governance must be tested as rigorously as business logic.
  \item Safety workflows demand deterministic transitions and explicit edge handling.
  \item Monorepo improves cohesion but requires dependency hygiene.
\end{itemize}

\section{Management Summary}
Project management outcomes show successful alignment between planning and technical delivery, while highlighting concrete maturity actions for production-level governance.

\chapter{Conclusion and Future Work}

\section{Global Conclusion}
This PFE delivered a complete SmartHome IoT platform integrating software architecture, IoT protocols, realtime systems, and role-based governance in a single coherent product. The final system exceeds a prototype dashboard and instead provides an operational platform with:
\begin{itemize}[leftmargin=1.2cm]
  \item Multi-role identity and onboarding model.
  \item House and unit lifecycle management.
  \item Device control over MQTT command channels.
  \item Resident permission governance with admin review.
  \item Safety-critical gas emergency monitoring and actuation loop.
  \item Realtime synchronization between backend events and frontend state.
\end{itemize}

\section{Objectives Evaluation}
The initial objectives defined in Chapter 1 are achieved at functional level. The platform demonstrates practical viability and can be used as foundation for industry-grade evolution.

\section{Academic Contribution}
From an academic perspective, the project demonstrates applied competence in:
\begin{itemize}[leftmargin=1.2cm]
  \item Requirements engineering and traceability.
  \item UML-driven architecture communication.
  \item Full-stack implementation and event-driven integration.
  \item Security-aware software design.
  \item Documentation and operational thinking.
\end{itemize}

\section{Industrial Readiness Assessment}
The solution is immediately suitable for demonstration, pilot deployments, and controlled production experiments. To reach enterprise-grade scale, the next stage should focus on distributed architecture decomposition, stronger observability, and compliance hardening.

\section{Future Work Roadmap}
\subsection{Technical Roadmap}
\begin{enumerate}[leftmargin=1.2cm]
  \item Introduce dedicated microservice for MQTT broker and telemetry ingestion.
  \item Add historical time-series analytics for consumption and anomaly detection.
  \item Implement policy engine for fine-grained permissions beyond role/action lists.
  \item Add mobile-native clients with push notification integration.
  \item Extend automation engine with rule editor and condition chaining.
\end{enumerate}

\subsection{Security Roadmap}
\begin{enumerate}[leftmargin=1.2cm]
  \item Add refresh-token strategy and session revocation list.
  \item Implement request rate limiting and anti-abuse controls.
  \item Apply MQTT ACL enforcement and topic-level authorization.
  \item Add secret rotation and formal key management process.
\end{enumerate}

\subsection{Operations Roadmap}
\begin{enumerate}[leftmargin=1.2cm]
  \item Add centralized logging with search and alerting.
  \item Integrate continuous deployment pipelines with staged gates.
  \item Build synthetic monitoring for critical user journeys.
\end{enumerate}

\section{Final Statement}
The SmartHome IoT Platform demonstrates how a thoughtful architecture can unify users, devices, and safety logic in one maintainable ecosystem. The project forms a robust baseline for both academic evaluation and real-world continuation.

%=============================================================================
%  Annexe A — Référence Complète de l'API
%=============================================================================
\chapter{Annexe A : Référence Complète de l'API}

\section{Vue d'Ensemble}

Tous les endpoints API utilisent \textbf{JSON} pour les corps de requête et de réponse.
Les endpoints protégés nécessitent un token \texttt{Bearer} dans l'en-tête
\texttt{Authorization}.

\begin{lstlisting}[style=codeblock,language=,caption={En-t\^{e}tes de Requ\^{e}te Communs}]
Content-Type: application/json
Authorization: Bearer <jwt_token>
\end{lstlisting}

\section{Endpoints d'Authentification (\api{/api/auth})}

\begin{longtable}{p{0.08\textwidth} p{0.32\textwidth} p{0.48\textwidth}}
\caption{API Authentification}\label{tab:api_auth}\\
\toprule
\textbf{Méthode} & \textbf{Chemin} & \textbf{Corps / Réponse} \\
\midrule
\endfirsthead
\toprule
\textbf{Méthode} & \textbf{Chemin} & \textbf{Corps / Réponse} \\
\midrule
\endhead
\bottomrule
\endfoot
POST & /api/auth/register/admin    & \texttt{\{nom,email,motdepasse\}} $\to$ \texttt{\{token,user,codeMaison\}} \\
POST & /api/auth/register/resident & \texttt{\{nom,email,mdp,codeMaison\}} $\to$ \texttt{\{token,user\}} \\
POST & /api/auth/register/agency   & \texttt{\{nom,email,mdp,codeAgence\}} $\to$ \texttt{\{token,user\}} \\
POST & /api/auth/login             & \texttt{\{email,motdepasse,totpToken?\}} $\to$ \texttt{\{token,user\}} \\
GET  & /api/auth/verify/:token     & Vérifie l'email, redirige vers le frontend \\
POST & /api/auth/forgot-password   & \texttt{\{email\}} $\to$ \texttt{\{message\}} \\
POST & /api/auth/reset-password    & \texttt{\{token,motdepasse\}} $\to$ \texttt{\{message\}} \\
GET  & /api/auth/me                & (JWT) $\to$ objet utilisateur complet \\
POST & /api/auth/resend-verification & \texttt{\{email\}} $\to$ \texttt{\{message\}} \\
\end{longtable}

\section{Endpoints Appareils (\api{/api/devices})}

\begin{longtable}{p{0.08\textwidth} p{0.32\textwidth} p{0.48\textwidth}}
\caption{API Appareils}\label{tab:api_devices}\\
\toprule
\textbf{Méthode} & \textbf{Chemin} & \textbf{Corps / Réponse} \\
\midrule
\endfirsthead
\toprule
\textbf{Méthode} & \textbf{Chemin} & \textbf{Corps / Réponse} \\
\midrule
\endhead
\bottomrule
\endfoot
GET    & /api/devices             & (Admin JWT) $\to$ liste \texttt{[Appareil]} de la maison \\
POST   & /api/devices             & \texttt{\{nom,type,topicMQTT,pièce,...\}} $\to$ \texttt{Appareil} \\
PUT    & /api/devices/:id         & Mise à jour partielle $\to$ \texttt{Appareil} mis à jour \\
DELETE & /api/devices/:id         & Supprime l'appareil, retourne \texttt{\{success:true\}} \\
GET    & /api/devices/:id         & Appareil unique par ID \\
POST   & /api/devices/:id/command & \texttt{\{command:\{state,brightness?\}\}} $\to$ \texttt{\{success\}} \\
\end{longtable}

\section{Endpoints Plan d'Étage (\api{/api/floorplan})}

\begin{longtable}{p{0.08\textwidth} p{0.32\textwidth} p{0.48\textwidth}}
\caption{API Plan d'Étage}\label{tab:api_floorplan}\\
\toprule
\textbf{Méthode} & \textbf{Chemin} & \textbf{Corps / Réponse} \\
\midrule
\endfirsthead
\toprule
\textbf{Méthode} & \textbf{Chemin} & \textbf{Corps / Réponse} \\
\midrule
\endhead
\bottomrule
\endfoot
GET  & /api/floorplan/state        & (JWT) $\to$ \texttt{\{appareils:[],houseState\}} \\
POST & /api/floorplan/command      & \texttt{\{deviceId,command\}} $\to$ \texttt{\{success\}} \\
POST & /api/floorplan/scene        & \texttt{\{scene:"Morning"|"Dinner"|"Cinematic"|"Sleep"\}} $\to$ \texttt{\{success,scene\}} \\
GET  & /api/floorplan/house-state  & (JWT) $\to$ document \texttt{HouseState} \\
POST & /api/floorplan/sync-defaults& (Admin) $\to$ synchronise les positions d'appareils \\
\end{longtable}

\section{Endpoints Sécurité Gaz (\api{/api/gas})}

\begin{longtable}{p{0.08\textwidth} p{0.32\textwidth} p{0.48\textwidth}}
\caption{API Gaz}\label{tab:api_gas}\\
\toprule
\textbf{Méthode} & \textbf{Chemin} & \textbf{Corps / Réponse} \\
\midrule
\endfirsthead
\toprule
\textbf{Méthode} & \textbf{Chemin} & \textbf{Corps / Réponse} \\
\midrule
\endhead
\bottomrule
\endfoot
GET  & /api/gas/state           & (JWT) $\to$ \texttt{\{niveau,seuil,vanneOuverte,urgence\}} \\
POST & /api/gas/clear-emergency & (Admin JWT) $\to$ \texttt{\{success:true\}} \\
POST & /api/gas/set-threshold   & (Admin JWT) \texttt{\{threshold:400\}} $\to$ \texttt{\{success\}} \\
\end{longtable}

\section{Endpoint Compagnon IA (\api{/api/companion})}

\begin{longtable}{p{0.08\textwidth} p{0.32\textwidth} p{0.48\textwidth}}
\caption{API Compagnon}\label{tab:api_companion}\\
\toprule
\textbf{Méthode} & \textbf{Chemin} & \textbf{Corps / Réponse} \\
\midrule
\endfirsthead
\toprule
\textbf{Méthode} & \textbf{Chemin} & \textbf{Corps / Réponse} \\
\midrule
\endhead
\bottomrule
\endfoot
POST & /api/companion/chat & \texttt{\{messages:[{role,content}]\}} $\to$ \texttt{\{reply,actionsExecuted\}} \\
\end{longtable}

\section{Endpoints Scénarios (\api{/api/scenarios})}

\begin{longtable}{p{0.08\textwidth} p{0.32\textwidth} p{0.48\textwidth}}
\caption{API Scénarios}\label{tab:api_scenarios}\\
\toprule
\textbf{Méthode} & \textbf{Chemin} & \textbf{Corps / Réponse} \\
\midrule
\endfirsthead
\toprule
\textbf{Méthode} & \textbf{Chemin} & \textbf{Corps / Réponse} \\
\midrule
\endhead
\bottomrule
\endfoot
GET    & /api/scenarios     & (Admin JWT) $\to$ liste des scénarios de la maison \\
POST   & /api/scenarios     & \texttt{\{nom,action,scheduleTime,isActive\}} $\to$ Scénario \\
PUT    & /api/scenarios/:id & Mise à jour partielle $\to$ Scénario mis à jour \\
DELETE & /api/scenarios/:id & Supprime le scénario $\to$ \texttt{\{success:true\}} \\
\end{longtable}

\section{Codes de Statut HTTP Courants}

\begin{center}
\begin{tabular}{cl}
\toprule
\textbf{Code} & \textbf{Signification dans cette API} \\
\midrule
200 & OK --- requête traitée avec succès \\
201 & Créé --- nouvelle ressource créée \\
400 & Mauvaise Requête --- échec de validation, doublon ou entrée invalide \\
401 & Non Autorisé --- token JWT manquant ou invalide \\
403 & Interdit --- authentifié mais rôle/permissions insuffisants \\
404 & Introuvable --- la ressource n'existe pas \\
429 & Trop de Requêtes --- limite de débit dépassée \\
500 & Erreur Interne Serveur --- défaillance inattendue du serveur \\
\bottomrule
\end{tabular}
\end{center}

\chapter{Appendix B: UML Source Catalog}

\section{Purpose}
This appendix provides text-source UML definitions (PlantUML style) so diagrams can be regenerated, versioned, and diffed as code artifacts.

\section{Use Case Diagram Source}
\begin{lstlisting}[style=codeblock]
@startuml
left to right direction
actor Agency
actor Admin
actor Resident

rectangle SmartHomePlatform {
  usecase "Manage Units" as UC1
  usecase "Review Owner Requests" as UC2
  usecase "Register Device" as UC3
  usecase "Activate Scenario" as UC4
  usecase "Review Permission Request" as UC5
  usecase "Request Permission" as UC6
  usecase "View Dashboard" as UC7
  usecase "Clear Gas Emergency" as UC8
}

Agency --> UC1
Agency --> UC2
Admin --> UC3
Admin --> UC4
Admin --> UC5
Admin --> UC8
Resident --> UC6
Resident --> UC7
Admin --> UC7
Agency --> UC7
@enduml
\end{lstlisting}

\section{Class Diagram Source (Core Domain)}
\begin{lstlisting}[style=codeblock]
@startuml
class User {
  +id: ObjectId
  +name: String
  +email: String
  +password: String
  +role: enum
  +houseCode: String
  +linkedAdmin: ObjectId
  +permissions: String[]
  +isEmailVerified: Boolean
}

class Device {
  +id: ObjectId
  +name: String
  +type: String
  +topic: String
  +room: String
  +status: String
  +state: Object
  +houseCode: String
}

class HouseState {
  +houseCode: String
  +lightingMode: String
  +lockdownMode: Boolean
  +awayMode: Boolean
  +gasLevel: Number
  +gasThreshold: Number
  +emergencyMode: Boolean
  +gasValveOpen: Boolean
}

class PermissionRequest {
  +requester: ObjectId
  +houseAdmin: ObjectId
  +houseCode: String
  +actionKey: String
  +actionLabel: String
  +status: enum
  +reviewer: ObjectId
}

class Unit {
  +unitCode: String
  +ownerID: ObjectId
  +ownerName: String
  +status: String
  +claimed: Boolean
  +isManaged: Boolean
}

class HouseOwnerRequest {
  +name: String
  +email: String
  +status: enum
  +accessCode: String
  +reviewedBy: ObjectId
}

User "1" --> "*" Device : owns
User "1" --> "*" PermissionRequest : requester
User "1" --> "*" PermissionRequest : reviewer/admin
Unit "1" --> "1" User : owner
HouseState "1" --> "*" Device : house scope
@enduml
\end{lstlisting}

\section{Sequence Source: Resident Permission Request}
\begin{lstlisting}[style=codeblock]
@startuml
actor ResidentUI
participant API as PermissionsAPI
database DB as MongoDB
participant Socket as SocketServer
participant AdminUI

ResidentUI -> PermissionsAPI: POST /permissions/request
PermissionsAPI -> DB: insert PermissionRequest
DB --> PermissionsAPI: created request
PermissionsAPI -> Socket: emit requestCreated
Socket -> AdminUI: requestCreated event
Socket -> ResidentUI: permissions:created event
@enduml
\end{lstlisting}

\section{Sequence Source: Admin Decision}
\begin{lstlisting}[style=codeblock]
@startuml
actor AdminUI
participant API as PermissionsAPI
database DB as MongoDB
participant Socket as SocketServer
participant ResidentUI

AdminUI -> PermissionsAPI: PATCH /permissions/admin/requests/:id
PermissionsAPI -> DB: update request status
alt approved
  PermissionsAPI -> DB: update resident permissions
end
PermissionsAPI -> Socket: emit requestStatusChanged
Socket -> ResidentUI: requestStatusChanged
Socket -> AdminUI: permissions:updated
@enduml
\end{lstlisting}

\section{Sequence Source: Gas Emergency}
\begin{lstlisting}[style=codeblock]
@startuml
participant Sensor
participant Broker as Aedes
participant Monitor as GasMonitor
database DB as HouseState
participant Socket
participant Dashboard

Sensor -> Broker: publish sensors/H1234/gas
Broker -> Monitor: publish event
Monitor -> DB: read/update gas state
alt level > threshold
  Monitor -> Broker: publish command/autogen/h1234/alarm
  Monitor -> Broker: publish command/autogen/h1234/gas-valve
end
Monitor -> Socket: emit house:gas-emergency
Socket -> Dashboard: alert payload
@enduml
\end{lstlisting}

\section{Activity Source: Onboarding Flow}
\begin{lstlisting}[style=codeblock]
@startuml
start
:User authenticates;
if (Has houseCode?) then (yes)
  :Redirect to /dashboard;
  stop
else (no)
  :Open /onboarding;
  if (Choose Join Unit?) then (yes)
    :POST /units/join;
    :Persist houseCode and isManaged;
  else (no)
    :POST /units/solo;
    :Persist houseCode and isSuperAdmin;
  endif
  :Redirect to /dashboard;
endif
stop
@enduml
\end{lstlisting}

\section{State Source: House Safety State Machine}
\begin{lstlisting}[style=codeblock]
@startuml
[*] --> Normal
Normal --> EmergencyPending : gasLevel > threshold
EmergencyPending --> EmergencyActive : trigger alarm + close valve
EmergencyActive --> Recovery : gasLevel <= threshold OR admin clear
Recovery --> Normal : emergencyMode=false
@enduml
\end{lstlisting}

\section{Component Source}
\begin{lstlisting}[style=codeblock]
@startuml
package Frontend {
  [React SPA]
  [AuthContext]
  [Dashboard Components]
}

package Backend {
  [Express API]
  [Socket.IO Service]
  [Aedes Broker]
  [Gas Monitor]
  [Route Modules]
  [Auth Middleware]
}

database MongoDB

[React SPA] --> [Express API] : REST
[React SPA] --> [Socket.IO Service] : realtime
[Aedes Broker] --> [Gas Monitor] : publish events
[Route Modules] --> MongoDB : persistence
[Gas Monitor] --> MongoDB : house state updates
[Gas Monitor] --> [Socket.IO Service] : alerts
@enduml
\end{lstlisting}

\section{Deployment Source}
\begin{lstlisting}[style=codeblock]
@startuml
node "Client Browser" {
  component "React SPA"
}

node "Node Host" {
  component "Express API"
  component "Socket.IO"
  component "Aedes MQTT Broker"
}

database "MongoDB"
node "IoT Network" {
  component "Sensors"
  component "Actuators"
}
cloud "Email Service"

"React SPA" --> "Express API" : HTTP
"React SPA" --> "Socket.IO" : WS
"Express API" --> "MongoDB" : Mongoose
"Sensors" --> "Aedes MQTT Broker" : MQTT publish
"Express API" --> "Aedes MQTT Broker" : command publish
"Express API" --> "Email Service" : SMTP/API
@enduml
\end{lstlisting}

\section{Sequence Source: Device Command}
\begin{lstlisting}[style=codeblock]
@startuml
actor UserUI
participant API as DeviceAPI
participant Broker as Aedes
participant Device

UserUI -> DeviceAPI: PUT /devices/:id/command
DeviceAPI -> Broker: publish command/<topic>
Broker -> Device: command payload
Device -> Broker: publish devices/<topic>
Broker -> DeviceAPI: publish event
@enduml
\end{lstlisting}

\section{DFD-style Logical Data Flow Source}
\begin{lstlisting}[style=codeblock]
@startuml
actor Resident
actor Admin
actor Agency

rectangle "Process P1: Auth and Access" as P1
rectangle "Process P2: Device and Scenario Control" as P2
rectangle "Process P3: Permission Governance" as P3
rectangle "Process P4: Gas Safety Monitor" as P4

database "D1 Users" as D1
database "D2 Devices" as D2
database "D3 HouseState" as D3
database "D4 PermissionRequests" as D4

Resident --> P1
Admin --> P1
Agency --> P1
P1 --> D1

Admin --> P2
P2 --> D2
P2 --> D3

Resident --> P3
Admin --> P3
P3 --> D4
P3 --> D1

P4 --> D3
P4 --> D2
@enduml
\end{lstlisting}

\section{Regeneration Instructions}
To regenerate these diagrams with PlantUML, save each block to a \texttt{.puml} file and run PlantUML locally or in CI. Export PNG or SVG and include in the report for publication versions.

\section{Appendix Summary}
Keeping UML sources as text enables architecture reviews, version control diffs, and future evolution tracking with minimal overhead.

\chapter{Appendix C: Installation Guide and Operations Runbook}

\section{Scope}
This appendix provides a complete practical guide to install, configure, run, troubleshoot, and maintain the SmartHome platform in development and pilot production environments.

\section{Prerequisites}
\subsection{Software}
\begin{itemize}[leftmargin=1.2cm]
  \item Node.js LTS (recommended v20+).
  \item npm (comes with Node.js).
  \item MongoDB server (local or managed).
  \item Optional: MQTT client tools for diagnostics.
  \item Optional: PM2 or container runtime for service supervision.
\end{itemize}

\subsection{System}
\begin{itemize}[leftmargin=1.2cm]
  \item Minimum 4 GB RAM for development.
  \item Open ports for backend HTTP and MQTT.
  \item Stable network for OAuth/email features.
\end{itemize}

\section{Repository Setup}
\begin{lstlisting}[style=codeblock,basicstyle=\ttfamily\small]
git clone <repository-url>
cd SmartHome-PFE
npm install
cd backend && npm install
cd ../frontend && npm install
\end{lstlisting}

\section{Environment Configuration}
Create backend environment file with keys such as:
\begin{lstlisting}[style=codeblock]
PORT=5000
MQTT_PORT=1883
MONGODB_URI=mongodb://127.0.0.1:27017/SmartHome
JWT_SECRET=replace_with_secure_secret
FRONTEND_URL=http://localhost:5173
BACKEND_URL=http://localhost:5000
ADMIN_ACCESS_CODE=ADMIN123
AGENCY_ACCESS_CODE=AGENCY2024
GOOGLE_CLIENT_ID=...
GOOGLE_CLIENT_SECRET=...
\end{lstlisting}

Frontend environment sample:
\begin{lstlisting}[style=codeblock]
VITE_API_BASE_URL=http://localhost:5000/api
\end{lstlisting}

\section{Local Run Commands}
\subsection{Monorepo Orchestrated Run}
\begin{lstlisting}[style=codeblock]
npm run start
\end{lstlisting}

\subsection{Separate Run}
\begin{lstlisting}[style=codeblock]
# terminal 1
cd backend
npm run dev

# terminal 2
cd frontend
npm run dev
\end{lstlisting}

\section{Health Verification Checklist}
\begin{enumerate}[leftmargin=1.2cm]
  \item Backend logs show successful MongoDB connection.
  \item HTTP API responds on configured port.
  \item Frontend starts and route guard logic works.
  \item MQTT TCP port accepts connections.
  \item Socket authenticated events delivered after login.
\end{enumerate}

\section{First-use Functional Smoke Tests}
\subsection{Auth Smoke}
\begin{itemize}[leftmargin=1.2cm]
  \item Register resident and login.
  \item Register or login admin.
  \item Confirm token-based API access.
\end{itemize}

\subsection{Device Smoke}
\begin{itemize}[leftmargin=1.2cm]
  \item Create a device as admin.
  \item Send a command to the device endpoint.
  \item Confirm MQTT publish and DB state update.
\end{itemize}

\subsection{Permission Smoke}
\begin{itemize}[leftmargin=1.2cm]
  \item Submit resident permission request.
  \item Review and approve from admin session.
  \item Confirm resident receives realtime status update.
\end{itemize}

\subsection{Gas Smoke}
\begin{itemize}[leftmargin=1.2cm]
  \item Publish gas reading below threshold.
  \item Publish reading above threshold.
  \item Confirm emergency flag and emitted alert.
\end{itemize}

\section{MQTT Diagnostic Recipes}
\subsection{Subscribe to command channel}
\begin{lstlisting}[style=codeblock]
mosquitto_sub -h localhost -p 1883 -t "command/#" -v
\end{lstlisting}

\subsection{Publish gas reading}
\begin{lstlisting}[style=codeblock]
mosquitto_pub -h localhost -p 1883 -t "sensors/H1234/gas" -m "{\"level\":650}"
\end{lstlisting}

\section{Common Incidents and Resolution}
\subsection{Incident: Cannot find module passport}
\textbf{Symptoms}: backend crashes at startup with module-not-found.

\textbf{Resolution}:
\begin{enumerate}[leftmargin=1.2cm]
  \item Verify backend dependencies installed.
  \item Install missing package in backend project.
  \item Restart backend runtime.
\end{enumerate}

\subsection{Incident: Socket authentication failed}
\textbf{Symptoms}: client cannot receive realtime updates.

\textbf{Resolution}:
\begin{enumerate}[leftmargin=1.2cm]
  \item Ensure valid JWT is sent in socket auth payload.
  \item Verify JWT secret consistency across environments.
  \item Confirm user still exists and token is not expired.
\end{enumerate}

\subsection{Incident: MQTT commands not reaching devices}
\textbf{Resolution actions}:
\begin{itemize}[leftmargin=1.2cm]
  \item Validate broker port and topic naming.
  \item Check payload serialization and parser logic.
  \item Confirm device simulator/client subscribes expected pattern.
\end{itemize}

\section{Security Hardening Checklist}
\begin{itemize}[leftmargin=1.2cm]
  \item Replace all default secrets and access codes.
  \item Enforce HTTPS and secure proxy headers.
  \item Restrict CORS origins.
  \item Add request rate limiting.
  \item Add audit logs for admin actions.
  \item Introduce secret rotation process.
\end{itemize}

\section{Backup and Restore Procedure}
\subsection{Backup}
\begin{lstlisting}[style=codeblock]
# example
mongodump --uri="mongodb://127.0.0.1:27017/SmartHome" --out=./backup
\end{lstlisting}

\subsection{Restore}
\begin{lstlisting}[style=codeblock]
# example
mongorestore --uri="mongodb://127.0.0.1:27017/SmartHome" ./backup/SmartHome
\end{lstlisting}

\section{Release Management}
\begin{enumerate}[leftmargin=1.2cm]
  \item Create release branch and freeze scope.
  \item Execute test suite and smoke checklist.
  \item Tag release candidate and deploy to staging.
  \item Validate monitoring and rollback readiness.
  \item Promote to production with changelog.
\end{enumerate}

\section{Rollback Strategy}
\begin{itemize}[leftmargin=1.2cm]
  \item Keep prior deploy artifacts accessible.
  \item Version database migrations and backups.
  \item Revert frontend static build and backend image atomically.
\end{itemize}

\section{Service Ownership Matrix}
\begin{table}[H]
\centering
\caption{Operational ownership examples}
\begin{tabular}{p{5cm}p{4cm}p{5cm}}
\toprule
Component & Primary Owner & Secondary Owner \\
\midrule
Backend API & Backend Engineer & Full-stack Engineer \\
Frontend SPA & Frontend Engineer & Full-stack Engineer \\
MongoDB & DevOps Engineer & Backend Engineer \\
MQTT Broker & IoT Engineer & Backend Engineer \\
Incident Response & DevOps Lead & Product Tech Lead \\
\bottomrule
\end{tabular}
\end{table}

\section{Capacity Planning Notes}
For scaling households and devices, monitor:
\begin{itemize}[leftmargin=1.2cm]
  \item peak concurrent socket connections,
  \item MQTT message rate,
  \item DB write amplification during burst events,
  \item API auth traffic under login spikes.
\end{itemize}

\section{Production Readiness Checklist}
\begin{itemize}[leftmargin=1.2cm]
  \item Secrets externalized and rotated.
  \item Monitoring and alerting active.
  \item Backup/restore tested.
  \item Security scan completed.
  \item Runbooks approved.
\end{itemize}

\section{Appendix Summary}
This runbook transforms the project from academic implementation to maintainable operational service by documenting reproducible procedures, controls, and recovery practices.

\chapter{Appendix D: Detailed Test Plan and Requirement Traceability}

\section{Objective}
This appendix formalizes validation strategy by mapping requirements to concrete test cases and acceptance criteria.

\section{Test Strategy Layers}
\begin{itemize}[leftmargin=1.2cm]
  \item Unit tests: isolated domain behavior and utilities.
  \item Integration tests: route + DB + realtime interactions.
  \item End-to-end tests: user journey validation by role.
  \item IoT simulation tests: MQTT telemetry and command loop.
\end{itemize}

\section{Environment Matrix}
\begin{longtable}{p{3.2cm}p{3.2cm}p{3cm}p{5.5cm}}
\toprule
Environment & Purpose & Data Type & Notes \\
\midrule
Local Dev & Feature development & Seed/demo & Fast iteration, debug logs \\
Integration & Contract validation & Synthetic realistic & Automated route-event tests \\
Staging & Pre-production checks & Sanitized close-to-prod & Deployment and performance smoke \\
Production & Live service & Real user/device data & Strict monitoring and controls \\
\bottomrule
\end{longtable}

\section{Requirement-to-Test Traceability Matrix}
\begin{longtable}{p{1.8cm}p{4cm}p{2.8cm}p{6.4cm}}
\toprule
Req ID & Requirement & Test ID & Validation Scope \\
\midrule
R-01 & Resident registration/login & T-AUTH-01..05 & Input validation, token issuance, role mapping \\
R-02 & Admin login with access code & T-AUTH-06..08 & Access code guard and role enforcement \\
R-03 & Unit join/solo onboarding & T-UNIT-01..06 & House assignment and capability flags \\
R-04 & Device registration and command & T-DEV-01..10 & CRUD scope and MQTT command publish \\
R-05 & Permission request workflow & T-PERM-01..12 & Lifecycle and realtime synchronization \\
R-06 & Gas emergency lifecycle & T-GAS-01..12 & Threshold transitions and actuator commands \\
R-07 & Agency unit governance & T-AGY-01..09 & Unit generation, update, analytics \\
R-08 & OAuth and account lifecycle & T-OAUTH-01..05 & Callback payload and auth consistency \\
\bottomrule
\end{longtable}

\section{Authentication Test Cases}
\subsection{T-AUTH-01 Resident Register Success}
\begin{itemize}[leftmargin=1.2cm]
  \item Preconditions: valid house code exists and invite policy satisfied.
  \item Steps: submit registration payload.
  \item Expected: HTTP 201, token present, role resident, houseCode persisted.
\end{itemize}

\subsection{T-AUTH-02 Resident Register Invalid House Code}
\begin{itemize}[leftmargin=1.2cm]
  \item Expected: HTTP 400 with invalid house code message.
\end{itemize}

\subsection{T-AUTH-03 Resident Login Unverified Email}
\begin{itemize}[leftmargin=1.2cm]
  \item Expected: HTTP 403 with EMAIL\_NOT\_VERIFIED code.
\end{itemize}

\subsection{T-AUTH-04 Admin Login Wrong Access Code}
\begin{itemize}[leftmargin=1.2cm]
  \item Expected: HTTP 401 and explicit invalid access code response.
\end{itemize}

\subsection{T-AUTH-05 Agency Login Wrong Role}
\begin{itemize}[leftmargin=1.2cm]
  \item Expected: HTTP 403 for non-agency account.
\end{itemize}

\section{Unit and Onboarding Test Cases}
\subsection{T-UNIT-01 Join Existing Unit Success}
\begin{itemize}[leftmargin=1.2cm]
  \item Verify user houseCode set to existing unitCode.
  \item Verify linked admin is assigned.
\end{itemize}

\subsection{T-UNIT-02 Join Unknown Unit}
\begin{itemize}[leftmargin=1.2cm]
  \item Expected: HTTP 404 no residence found.
\end{itemize}

\subsection{T-UNIT-03 Create Solo Unit}
\begin{itemize}[leftmargin=1.2cm]
  \item Expected: unique code generated, isManaged false, isSuperAdmin true.
\end{itemize}

\section{Device Test Cases}
\subsection{T-DEV-01 Register Device as Admin}
\begin{itemize}[leftmargin=1.2cm]
  \item Expected: device created with house-scoped ownership.
\end{itemize}

\subsection{T-DEV-02 Register Device Duplicate Topic}
\begin{itemize}[leftmargin=1.2cm]
  \item Expected: HTTP 400 duplicate topic message.
\end{itemize}

\subsection{T-DEV-03 Send Device Command}
\begin{itemize}[leftmargin=1.2cm]
  \item Expected: MQTT publish to command/<topic>, DB state update.
\end{itemize}

\subsection{T-DEV-04 Delete Device Unauthorized}
\begin{itemize}[leftmargin=1.2cm]
  \item Expected: role guard blocks deletion.
\end{itemize}

\section{Permission Workflow Test Cases}
\subsection{T-PERM-01 Resident Creates Request}
\begin{itemize}[leftmargin=1.2cm]
  \item Expected: pending request created, admin notification emitted.
\end{itemize}

\subsection{T-PERM-02 Prevent Duplicate Pending Request}
\begin{itemize}[leftmargin=1.2cm]
  \item Expected: HTTP 409 conflict.
\end{itemize}

\subsection{T-PERM-03 Admin Approves Request}
\begin{itemize}[leftmargin=1.2cm]
  \item Expected: status approved, reviewer fields set, resident permissions merged.
\end{itemize}

\subsection{T-PERM-04 Admin Denies Request}
\begin{itemize}[leftmargin=1.2cm]
  \item Expected: status denied, no permission merge.
\end{itemize}

\subsection{T-PERM-05 Read-status Synchronization}
\begin{itemize}[leftmargin=1.2cm]
  \item Expected: read timestamps update for relevant actor side.
\end{itemize}

\section{Gas Safety Test Cases}
\subsection{T-GAS-01 Reading Below Threshold}
\begin{itemize}[leftmargin=1.2cm]
  \item Expected: emergency false, state updated with latest gasLevel.
\end{itemize}

\subsection{T-GAS-02 Reading Above Threshold}
\begin{itemize}[leftmargin=1.2cm]
  \item Expected: emergency true, alarm publish, valve close publish.
\end{itemize}

\subsection{T-GAS-03 Emergency Clear by Admin}
\begin{itemize}[leftmargin=1.2cm]
  \item Expected: emergency false and reset event broadcast.
\end{itemize}

\subsection{T-GAS-04 Threshold Update Bounds}
\begin{itemize}[leftmargin=1.2cm]
  \item Expected: threshold clamped to allowed range and persisted.
\end{itemize}

\section{Agency and Owner-request Test Cases}
\subsection{T-AGY-01 Generate Unit}
\begin{itemize}[leftmargin=1.2cm]
  \item Expected: unique unclaimed unit record created.
\end{itemize}

\subsection{T-AGY-02 Approve Owner Request}
\begin{itemize}[leftmargin=1.2cm]
  \item Expected: status approved, access code generated, email dispatch invoked.
\end{itemize}

\subsection{T-AGY-03 Reject Owner Request}
\begin{itemize}[leftmargin=1.2cm]
  \item Expected: status rejected with optional note.
\end{itemize}

\section{Socket Event Validation Cases}
\begin{longtable}{p{3cm}p{4cm}p{7cm}}
\toprule
Test ID & Event & Expected Result \\
\midrule
T-SOCK-01 & requestCreated & Delivered to house-admin room and admin user room \\
T-SOCK-02 & requestStatusChanged & Delivered to requester and reviewing admin \\
T-SOCK-03 & permissionsUpdated & Delivered to requester user room with new permissions \\
T-SOCK-04 & house:gas-emergency & Delivered to house room and agency global in emergency \\
T-SOCK-05 & dashboard:state-updated & Delivered to house room after scene activation \\
\bottomrule
\end{longtable}

\section{Performance and Load Test Ideas}
\begin{itemize}[leftmargin=1.2cm]
  \item Simulate N concurrent socket clients per house.
  \item Burst MQTT telemetry and observe update latency.
  \item Run API endpoint load profile for auth and permissions.
\end{itemize}

\section{Quality Gates for CI}
\begin{enumerate}[leftmargin=1.2cm]
  \item Lint success in backend and frontend.
  \item Unit tests pass threshold.
  \item Integration smoke suite green.
  \item Security static checks with zero critical findings.
\end{enumerate}

\section{Defect Lifecycle Template}
\begin{longtable}{p{2.5cm}p{3cm}p{3cm}p{5.5cm}}
\toprule
Field & Example & Owner & Notes \\
\midrule
Defect ID & BUG-047 & QA Engineer & Unique tracking identifier \\
Severity & High & Tech Lead & Impact to safety/auth/business \\
Status & In Progress & Developer & Open, In Progress, Ready, Closed \\
Root Cause & Missing role guard & Backend Engineer & Include preventive action \\
\bottomrule
\end{longtable}

\section{Exit Criteria for Release}
Release candidate can be accepted when:
\begin{itemize}[leftmargin=1.2cm]
  \item no blocker defects remain,
  \item all critical test cases pass,
  \item security checks pass baseline,
  \item smoke tests succeed in staging.
\end{itemize}

\section{Appendix Summary}
This appendix turns requirements into measurable validation assets, enabling reproducible quality governance over future project iterations.


%=============================================================================
%  BIBLIOGRAPHY
%=============================================================================
\bibliographystyle{plain}
\bibliography{references}

\end{document}
